% =============================================================================
% Appendix - Mathematical Formulations
% =============================================================================
% Collects every display formula in the thesis, with its symbols, mechanism,
% source attribution and implementation status.
%
% INCLUSION: this file is an APPENDIX, not a chapter. In main.tex:
%
%     \appendix
%     % =============================================================================
% Appendix - Mathematical Formulations
% =============================================================================
% Collects every display formula in the thesis, with its symbols, mechanism,
% source attribution and implementation status.
%
% INCLUSION: this file is an APPENDIX, not a chapter. In main.tex:
%
%     \appendix
%     % =============================================================================
% Appendix - Mathematical Formulations
% =============================================================================
% Collects every display formula in the thesis, with its symbols, mechanism,
% source attribution and implementation status.
%
% INCLUSION: this file is an APPENDIX, not a chapter. In main.tex:
%
%     \appendix
%     % =============================================================================
% Appendix - Mathematical Formulations
% =============================================================================
% Collects every display formula in the thesis, with its symbols, mechanism,
% source attribution and implementation status.
%
% INCLUSION: this file is an APPENDIX, not a chapter. In main.tex:
%
%     \appendix
%     \include{contents/appendix_formulas}
%
% It must come after \appendix. Including it as a numbered chapter would shift
% the chapter numbering that chapter3_shortened.tex's divergence ledger depends
% on (it cites ten section numbers literally: 3.1.2, 3.1.3, 3.2.6, 3.4.3,
% 3.4.7, 3.5.7, 3.6.6, 3.7.7, 3.8.6, 3.9.7).
%
% CITATION POLICY: this file adds NO new bibliography entries. Works cited only
% inside the reviewed papers (Vaswani, Lin et al., Ioffe & Szegedy, Wu et al.)
% are attributed in prose to the corpus paper that reports them, per the policy
% stated in the header of bib/thesis.bib. Adding an entry would renumber the
% whole bibliography, since tools/bibliography.py numbers by first appearance.
%
% All labels in this file are prefixed eq: and collide with nothing.
% =============================================================================

\chapter{Mathematical Formulations}
\label{ch:formulas}

This appendix collects every display formula that appears in the body of the
thesis, in one place, with a consistent treatment: the expression as it is
written in its own chapter, the section it belongs to, the meaning of each
symbol, the mechanism it implements, the work it comes from, and whether this
thesis implements it as written.

\section*{How to read this appendix}
\addcontentsline{toc}{section}{How to read this appendix}

\textbf{Ordering.} Sequential by chapter, in the order the mathematics is
introduced: Chapter~3 (\autoref{sec:eq-ch3}), Chapter~4
(\autoref{sec:eq-ch4}), then Chapter~2 (\autoref{sec:eq-ch2}), which carries
the bulk of the derivations. Chapters~1, 5 and~6 contain no display
mathematics.

\textbf{Reproduction.} Every expression is reproduced exactly as it is
typeset in its own chapter --- same symbols, same subscripts, same
conventions. Where two chapters state the same relation in different notation,
both forms are given inside the one entry rather than being silently
reconciled.

\textbf{Each formula appears once.} Chapter~2 develops the mechanism and
Chapter~3 reports the literature's version and any divergence from it, so
several relations occur in both. Those are presented once, at the location
listed first in the ordering above, with a cross-reference recording every
section that states them. No entry is duplicated.

\textbf{Attribution.} Citations use only keys already present in
\texttt{bib/thesis.bib}. Several formulas originate in works that the
bibliography deliberately omits, because this thesis's corpus is the set of
reviewed micro-expression papers rather than the full ancestry of every
technique they use. Those are attributed in prose to the reviewed paper that
brings them to this task --- focal loss through Zhao et
al.~\cite{zhaoS2021}, self-attention and the transformer through Dosovitskiy
et al.~\cite{dosovitskiy2021} and Zhang et al.~\cite{zhang2022}, Eulerian
magnification through Bai et al.~\cite{bai2021}. Formulas that are standard
textbook material with no corpus source, such as batch normalisation and the
convolution parameter count, carry no citation and are marked as such.

\textbf{Implementation status.} Each entry closes by stating whether the
thesis implements the formula as written. Three do not, and those are declared
divergences recorded in the divergence ledger of \autoref{tab:divergence-ledger}
rather than defects to be corrected here.

\begin{table}[htbp]
	\centering
	\footnotesize
	\begin{tabular}{@{}l l l l@{}}
		\hline
		\textbf{Entry} & \textbf{Quantity} & \textbf{Attributed to} & \textbf{As written?} \\
		\hline
		\multicolumn{4}{@{}l}{\textit{Chapter 3 --- Literature Review}} \\
		\autoref{eq:uf1}           & Per-class F1 and UF1          & \cite{see2019}                        & yes \\
		\autoref{eq:strain-frobenius} & Four-term strain magnitude & \cite{liong2014a,liong2014b,liong2016} & \textbf{no} \\
		\autoref{eq:strain-implemented} & Three-term strain magnitude & declared divergence                & --- \\
		\autoref{eq:simam}         & SimAM energy and gate         & \cite{yang2021}                       & partly \\
		\autoref{eq:focal}         & Focal loss                    & \cite{zhaoS2021}                      & partly \\
		\multicolumn{4}{@{}l}{\textit{Chapter 4 --- Methodology}} \\
		\autoref{eq:posenc}        & Sinusoidal positional encoding & no corpus source                     & yes \\
		\autoref{eq:fold-ceiling}  & Mean-of-folds macro F1 ceiling & this thesis                          & n/a \\
		\multicolumn{4}{@{}l}{\textit{Chapter 2 --- Background}} \\
		\autoref{eq:evm-taylor}    & Small-motion approximation    & \cite{bai2021}                        & yes \\
		\autoref{eq:evm-amplified} & Amplified signal              & \cite{bai2021}                        & \textbf{no} \\
		\autoref{eq:brightness}    & Brightness constancy          & \cite{liong2019a}                     & yes \\
		\autoref{eq:flow-constraint} & Optical flow constraint     & \cite{liong2019a}                     & yes \\
		\autoref{eq:strain-tensor} & Symmetric displacement gradient & \cite{shreve2011}                   & yes \\
		\autoref{eq:conv-params}   & Convolution parameter count   & standard                              & yes \\
		\autoref{eq:batchnorm}     & Batch normalisation           & standard, no corpus source            & yes \\
		\autoref{eq:attention}     & Scaled dot-product attention  & via \cite{dosovitskiy2021}            & yes \\
		\autoref{eq:label-smoothing} & Label smoothing             & named by \cite{dosovitskiy2021}       & yes \\
		\autoref{eq:prf}           & Precision, recall, F1         & standard                              & yes \\
		\autoref{eq:macro-f1}      & Macro-averaged F1             & standard; \cite{see2019} pooled       & yes \\
		\hline
	\end{tabular}
	\caption[Index of formulations]{Index of the formulations collected in this
	appendix, with the work each is attributed to and whether this thesis
	implements it as written.}
	\label{tab:formula-index}
\end{table}


% =============================================================================
\section{Chapter 3 --- Literature Review}
\label{sec:eq-ch3}
% =============================================================================

% -----------------------------------------------------------------------------
\subsection{Per-class F1 and Unweighted F1}
\label{eq:uf1}
% -----------------------------------------------------------------------------

\textbf{Appears at} \autoref{sec:megc-metrics}.

\[
	F1_c = \frac{2 \cdot TP_c}{2 \cdot TP_c + FP_c + FN_c}, \qquad \text{UF1} = \frac{1}{C}\sum_c F1_c
\]

\textbf{Symbols.} $TP_c$, $FP_c$ and $FN_c$ are the true positives, false
positives and false negatives for class $c$; $C$ is the number of classes,
three here; $F1_c$ is the harmonic mean of precision and recall for class $c$,
written in the algebraically equivalent form that avoids computing either
separately.

\textbf{What it does.} UF1 scores a classifier without letting a frequent
class dominate, by computing a score per class and averaging the scores with
equal weight rather than averaging over examples. The construction that
matters is not visible in the algebra: the counts $TP_c$, $FP_c$ and $FN_c$
are \emph{accumulated across all $k$ folds before any $F1_c$ is computed}.
That pooling is what distinguishes UF1 from the average of per-fold macro F1
scores, which is a different estimator with a different expectation
(\autoref{sec:pooled-vs-fold}, and \autoref{eq:fold-ceiling} for what it costs).

\textbf{Source.} The metric convention of the Second Facial Micro-Expression
Grand Challenge, See et al.~\cite{see2019}, which mandates UF1 and Unweighted
Average Recall in place of plain accuracy on the explicit grounds of class
imbalance.

\textbf{As implemented.} Yes, and it is the thesis's primary metric. Pooled
macro F1 as computed here is identical in construction to MEGC's UF1. The
thesis adopts the challenge's metric definitions while declining its composite
training regime.

% -----------------------------------------------------------------------------
\subsection{Four-term strain magnitude (the published convention)}
\label{eq:strain-frobenius}
% -----------------------------------------------------------------------------

\textbf{Appears at} \autoref{sec:strain-tensor}.

\[
    \varepsilon_{\mathrm{F}}
    =\sqrt{\varepsilon_{xx}^{2}+\varepsilon_{yy}^{2}
        +\varepsilon_{xy}^{2}+\varepsilon_{yx}^{2}}
    =\sqrt{\varepsilon_{xx}^{2}+\varepsilon_{yy}^{2}
        +2\varepsilon_{xy}^{2}}.
\]

\textbf{Symbols.} $\varepsilon_{xx}$ and $\varepsilon_{yy}$ are the normal
components of the symmetric displacement-gradient tensor, describing local
stretching or compression along each axis; $\varepsilon_{xy}=\varepsilon_{yx}$
are its shear components. The tensor itself is defined at
\autoref{eq:strain-tensor}. The second equality follows from the tensor's
symmetry, which makes the two off-diagonal entries equal.

\textbf{What it does.} It reduces the four-component tensor at each pixel to
one non-negative scalar, the Frobenius norm, describing deformation
\emph{intensity} while discarding its sign and direction. Because the tensor is
symmetric, the norm counts shear twice --- once for each off-diagonal entry ---
which is what produces the factor of two.

\textbf{Source.} The Liong strain series~\cite{liong2014a,liong2014b,liong2016}.
Shreve et al.~\cite{shreve2011} establish the tensor; the Liong papers use this
four-component magnitude as the descriptor.

\textbf{As implemented.} \textbf{No.} The implementation computes the
three-term form of \autoref{eq:strain-implemented} instead. This is the
thesis's most consequential divergence and is treated in the next entry.

% -----------------------------------------------------------------------------
\subsection{Three-term strain magnitude (the implemented form)}
\label{eq:strain-implemented}
% -----------------------------------------------------------------------------

\textbf{Appears at} \autoref{sec:strain-divergence}, and again as
$\varepsilon_{\mathrm{mag}}$ at \autoref{sec:optical-strain} in Chapter~2.

\[
    \varepsilon_{\mathrm{os}}
    =\sqrt{\varepsilon_{xx}^{2}+\varepsilon_{yy}^{2}
        +\varepsilon_{xy}^{2}},
    \qquad
    \varepsilon_{xy}=\tfrac{1}{2}
        \left(\frac{\partial u}{\partial y}
             +\frac{\partial v}{\partial x}\right).
\]

Chapter~2 writes the same scalar as

\[
    \varepsilon_{\mathrm{mag}}
       =\sqrt{\varepsilon_{xx}^2+\varepsilon_{yy}^2
                    +\varepsilon_{xy}^2}.
\]

\textbf{Symbols.} As \autoref{eq:strain-frobenius}, with $u$ and $v$ the
horizontal and vertical components of the estimated flow field. In the
discrete pipeline $u$ and $v$ are displacements in pixels between selected
frames, not physical velocities.

\textbf{What it does.} The same reduction to a deformation-intensity scalar,
but counting the shear contribution once rather than twice. The consequence is
a change in relative weighting, not a removable constant: the four-term form
weights pure shear by $\sqrt{2}$ relative to this one. Independent min--max
normalisation removes a uniform scale factor but does not generally remove
this difference, because normal and shear contributions vary across pixels and
across time.

\textbf{Source.} \textbf{None --- this is a declared divergence.} STSTNet
\cite{liong2019b} cannot establish a precedent for the three-term form,
because its printed magnitude equation is typographically defective in two
independent ways. It repeats $\partial u$ in the mixed derivative where the
symmetric tensor calls for $\partial v / \partial x$ alongside $\partial u /
\partial y$, and it places the factor $\tfrac{1}{2}$ \emph{outside} the
squared derivative sum. Correcting the repeated derivative while retaining
that placement gives

\[
    \tfrac{1}{2}
       \left(\frac{\partial u}{\partial y}
            +\frac{\partial v}{\partial x}\right)^2
       =2\varepsilon_{xy}^{2},
\]

which agrees with the \emph{four}-term convention, not the three-term one.
Moving the factor inside the square would instead yield the implemented form,
but that is a further change to the printed expression. The intended
implementation cannot be settled from the published equation alone, so the
thesis declares a departure from the unambiguous four-term convention rather
than claiming to reproduce STSTNet's magnitude.

\textbf{As implemented.} This is what the code computes. Spatial derivatives
come from NumPy's \texttt{gradient} at one-pixel spacing, using central
differences in the array interior; the spacing follows Liong et
al.~\cite{liong2014a} even though the scalar reduction does not. The two
conventions were never compared experimentally here, so their practical
equivalence is not assumed. This is row 4 of the divergence ledger
(\autoref{tab:divergence-ledger}). \textbf{Neither expression should be
``corrected'' --- the difference between them is the finding.}

% -----------------------------------------------------------------------------
\subsection{SimAM: energy and gate}
\label{eq:simam}
% -----------------------------------------------------------------------------

\textbf{Appears at} \autoref{sec:simam-energy}, and at \autoref{sec:simam} in
Chapter~2. Chapter~3 writes it in terms of the intermediate $d_i$ and $s^2$:

\[
    d_i=(x_i-\mu)^2,\qquad
    s^2=\frac{1}{M-1}\sum_{i=1}^{M}d_i,\qquad
    E_i^{-1}=\frac{d_i}{4(s^2+\lambda)}+\frac{1}{2},
\]
\[
    \widetilde{x}_i=x_i\,\operatorname{sigmoid}(E_i^{-1}),
    \qquad \mu=\frac{1}{M}\sum_{i=1}^{M}x_i.
\]

Chapter~2 states the same gate with the variance written as $v$:

\[
    \mu=\frac1M\sum_i x_i,\qquad
    v=\frac1{M-1}\sum_i(x_i-\mu)^2,
\]
\[
    E_i^{-1}=\frac{(x_i-\mu)^2}{4(v+\lambda)}+\frac12,
    \qquad \widetilde x_i=x_i\operatorname{sigmoid}(E_i^{-1}).
\]

\textbf{Symbols.} $x_i$ is one activation; $\mu$ and $s^2$ (equivalently $v$)
are the mean and variance over the statistical context shared by that
channel; $M$ is the number of positions in that context; $\lambda$ is a
positive regularisation constant, fixed at $10^{-4}$ here; $E_i^{-1}$ is the
reciprocal energy, the importance score; $\widetilde{x}_i$ is the refined
activation.

\textbf{What it does.} It implements a distinctiveness prior: an activation
that differs from its channel context is scored as more important, and the
score multiplicatively rescales the feature map rather than being added to it.
The score is the closed-form minimiser of a regularised separation energy, so
the $+\tfrac{1}{2}$ follows from the reciprocal-energy expression rather than
being an arbitrary offset. The sigmoid bounds the multiplicative weight while
preserving the ordering of the scores; it does not produce a hard mask of
inactive regions. No weight or bias is learned anywhere in the module, so it
contributes \textbf{zero learnable parameters} regardless of channel count.

Note what the prior is and is not. It is a statement about activation
statistics, not a detector of facial action units or of minority classes. A
distinctive response may encode useful deformation, or noise, or an artefact.

\textbf{Source.} Yang et al.~\cite{yang2021}.

\textbf{As implemented.} Partly --- with two documented distinctions, both
deliberate.

\begin{enumerate}
	\item \textbf{The variance denominator.} The shared variance printed with
	Yang et al.'s equation~(5) divides by $M$; the reference code in their
	Figure~3 divides by $M-1$. The forms above follow the \emph{code}, as does
	this thesis, which applies Bessel's correction over $n = D \cdot H \cdot W
	- 1$ (\autoref{sec:simam-placement}).
	\item \textbf{The statistical context.} Here $M = THW$ --- whole-clip
	statistics within each channel --- against the original image
	formulation's $M = HW$ per frame. The original module produces a separate
	weight at each channel--height--width position, so its ``three-dimensional
	weights'' describe $C \times H \times W$ and not a temporal volume.
	Extending the context across frames is an additional design decision, not
	part of the source formulation, and is row 10 of the divergence ledger.
	Whole-clip statistics do not guarantee that apex frames receive the
	highest weights.
\end{enumerate}

The value $\lambda = 10^{-4}$ is likewise inherited rather than derived: Yang
et al. sweep $\lambda$ from $10^{-1}$ to $10^{-6}$ and select $10^{-4}$ for
CIFAR, repeating the selection separately for ImageNet and choosing $0.1$. It
is a dataset-selected operating point, not a universal setting justified by
the formula.

% -----------------------------------------------------------------------------
\subsection{Focal loss}
\label{eq:focal}
% -----------------------------------------------------------------------------

\textbf{Appears at} \autoref{sec:imbalance-focal}, and at
\autoref{sec:focal-loss} in Chapter~2.

\[
	\text{FL}(p_t) = -\alpha_t (1 - p_t)^{\gamma}\log p_t
\]

\textbf{Symbols.} $p_t$ is the model's predicted probability for the true
class; $\gamma$ is the focusing exponent, described by Zhao et al. as ``the
balance factor for loss''; $\alpha_t$ is a per-class multiplier, ``the weight
balance factor for samples''. With $\alpha_t = 1$ and $\gamma = 0$ the
expression reduces to ordinary cross-entropy, $-\log p_t$.

\textbf{What it does.} Two independent corrections, which are easy to
conflate and which this thesis keeps separate deliberately. The modulating
factor $(1-p_t)^{\gamma}$ re-weights by \textbf{difficulty}: for $\gamma > 0$
it shrinks the contribution of confident, correct examples relative to hard
ones. It does not inspect class frequency at all, so a difficult example from
any class retains a large contribution --- though under skew the examples it
up-weights are disproportionately from the minority classes. The separate
$\alpha_t$ multiplier re-weights by \textbf{class frequency}. Zhao et al.
treat $\gamma$ as a constant and $\alpha$ as dataset-specific, to be set by
cross-validation.

\textbf{Source.} Zhao et al.~\cite{zhaoS2021}, who bring focal loss into
micro-expression recognition ``to alleviate the inefficient model training
caused by the class imbalance in the MEs Datasets''. They adopt the
formulation developed for one-stage object detection, where the imbalance
being corrected is between foreground and background regions rather than
between emotion classes. Per the bibliography's stated policy, the original
detection paper carries no entry here and the formulation is attributed to the
reviewed paper that reports it.

\textbf{As implemented.} Partly. The difficulty term is used as written, with
$\gamma = 2$ fixed throughout, matching Zhao et al.'s reported exponent --- an
inherited setting, not a measured optimum for this corpus. The implementation
computes $p_t$ from the hard target label using log-softmax, then multiplies a
\emph{label-smoothed} cross-entropy term (\autoref{eq:label-smoothing}) by
$(1-p_t)^{\gamma}$, with the class weight applied afterwards and optionally.
The frequency term $\alpha_t$ is where this thesis departs: frequency
correction is assigned to the sampler instead, and when balanced sampling is
active a runtime guard disables the loss-side class weights, leaving focal
difficulty weighting in place. Conflating the two factors is precisely the
error recorded at \autoref{sec:imbalance-double-correction}.

% -----------------------------------------------------------------------------
\subsection{Inline quantities}
\label{eq:ch3-inline}
% -----------------------------------------------------------------------------

Three quantities in Chapter~3 are not set as display mathematics but carry the
same evidential weight.

\textbf{Flow magnitude, $\sqrt{u^2+v^2}$} (\autoref{sec:brightness-constancy}).
The scalar magnitude of the flow field, described by Zhao et
al.~\cite{zhaoS2021} alongside the component fields $u(x,y)$ and $v(x,y)$.
This thesis retains the two components as separate input channels rather than
reducing them to this magnitude, so the direction of motion is preserved.

\textbf{Effective frame rate, $6600/N$ fps} (\autoref{sec:tim-implications}).
Sampling 33 frames from a clip of $N$ frames recorded at 200~fps yields this
effective rate: \textbf{100~fps} at the median clip length $N=66$, which is
SMIC's recording rate and half what CASME~II was built to provide, and roughly
\textbf{52~fps} for the longest clip at $N=126$, below the original CASME
database's 60~fps~\cite{yan2014}. Eighty-one of the 156 clips are downsampled
by a factor of two or more. No reviewed paper reports this cost.

\textbf{True inter-frame interval, $N/(200 \times 32)$ seconds}
(\autoref{sec:tim-implications}). Magnification runs on the 33 sampled frames
at the nominal 200~fps, but those frames span the clip's full duration, so the
real interval between them is this rather than $1/200$. A band specified as
5--25~Hz therefore corresponds to roughly 2.5--12.5~Hz physically at the
median, and the discrepancy grows with clip length. The realised pass-band is
consequently clip-dependent --- neither the 5--25~Hz derived from the duration
rule of Bai et al.~\cite{bai2021} nor constant across the corpus. This is a
declared divergence; magnifying before subsampling is declared as further
work.


% =============================================================================
\section{Chapter 4 --- Methodology}
\label{sec:eq-ch4}
% =============================================================================

% -----------------------------------------------------------------------------
\subsection{Sinusoidal positional encoding}
\label{eq:posenc}
% -----------------------------------------------------------------------------

\textbf{Appears at} \autoref{sec:transformer-config}.

\[
	\mathrm{PE}(p, 2i) = \sin\!\left(\frac{p}{10000^{2i/96}}\right), \qquad \mathrm{PE}(p, 2i+1) = \cos\!\left(\frac{p}{10000^{2i/96}}\right).
\]

\textbf{Symbols.} $p$ is the position in the sequence, one of the 32 time
steps; $i$ indexes feature pairs within the embedding, so $2i$ and $2i+1$ are
the even and odd feature positions; 96 is $d_{\text{model}}$, the embedding
width, appearing here as a literal rather than a symbol because the
configuration fixes it.

\textbf{What it does.} It supplies the order information that self-attention
cannot recover on its own. Without it, attention is permutation-equivariant:
permuting the input tokens permutes the outputs correspondingly, so a shuffled
clip is indistinguishable from an ordered one (\autoref{sec:self-attention}).
Each feature pair oscillates at a different frequency, geometrically spaced by
the $10000^{2i/96}$ denominator, so the vector at each position is distinct
and positions at a fixed offset stand in a consistent linear relation.

\textbf{Source.} No corpus source. The encoding is Vaswani et al.'s, from the
original transformer, which the bibliography deliberately omits under the
policy described at the head of this appendix; Chapter~4 states it as the
standard form without attribution. The thesis reaches the architecture through
Zhang et al.'s SLSTT~\cite{zhang2022} and its Vision Transformer
foundation~\cite{dosovitskiy2021}, as traced at
\autoref{sec:transformer-review}. Note that ViT itself uses \emph{learned}
positional embeddings, not these fixed sinusoids.

\textbf{As implemented.} Yes, as written. It is a fixed buffer, not a learned
parameter, and contributes nothing to any parameter count. Two details of the
surrounding configuration are easily missed: the positional-encoding stage
applies its own \texttt{Dropout(p=0.1)} to the sequence \emph{after} the
encoding is added, which is a separate application from the encoder layers'
dropout at the same rate; and although the sinusoids could be evaluated at
arbitrary positions, the implementation uses a finite precomputed buffer, so
generalisation to unseen sequence lengths is neither exercised nor guaranteed.

% -----------------------------------------------------------------------------
\subsection{The mean-of-folds macro F1 ceiling}
\label{eq:fold-ceiling}
% -----------------------------------------------------------------------------

\textbf{Appears at} \autoref{sec:mean-of-folds-rejected}.

\[
	\frac{10 \cdot \tfrac13 + 8 \cdot \tfrac23 + 7 \cdot 1}{25} = \frac{3.333\ldots + 5.333\ldots + 7}{25} \approx 0.627.
\]

\textbf{Symbols.} The three numerator terms are the 25 leave-one-subject-out
folds grouped by how many of the three classes appear in each fold's held-out
ground truth: \textbf{10} folds contain one class, \textbf{8} contain two,
\textbf{7} contain all three. The fractions $\tfrac13$, $\tfrac23$ and $1$ are
the maximum macro F1 each group can attain. The denominator is the fold count.

\textbf{What it does.} It computes an upper bound, fixed by protocol
structure before any model is trained. A macro F1 computed \emph{within} a
fold is the unweighted mean of that fold's three per-class F1 scores, and a
class absent from that fold's ground truth contributes zero no matter what the
model predicts. A single-class fold is therefore capped at $1/3$ and a
two-class fold at $2/3$. Averaging across folds averages seventeen sub-unity
ceilings into every score, bounding the mean-of-folds macro F1 at
approximately \textbf{0.627} --- a property of which subjects were held out,
identical for all twelve configurations.

Pooling avoids the ceiling entirely: accumulating $TP$, $FP$ and $FN$ across
all 25 folds and all 156 clips before computing any per-class F1 means a class
locally absent from one fold no longer forces a zero into its contribution.
That is the construction of \autoref{eq:uf1}.

\textbf{Source.} This thesis's own arithmetic over the fold composition of
\autoref{fig:fold-composition}. It is not a published result.

\textbf{As implemented.} Not applicable --- this is an argument for rejecting
a metric, not a component. Its consequence is directly implemented: pooled
macro F1 is the primary reported metric, and the mean-of-folds quantity stored
under \texttt{macro\_f1} in the run artefacts is not. A metric with a hard
ceiling below 0.7 cannot serve as a headline result. The same fold composition
biases mean-of-folds \emph{accuracy} in the same direction without imposing a
hard ceiling; for the full configuration it exceeds pooled accuracy by
\textbf{6.3 points}.

% -----------------------------------------------------------------------------
\subsection{Inline quantities}
\label{eq:ch4-inline}
% -----------------------------------------------------------------------------

\textbf{Sampler weight, $1/n_c$} (\autoref{sec:sampler-config}). Every
training example is assigned this weight by the
\texttt{WeightedRandomSampler}, where $n_c$ is the number of training examples
of that example's class \emph{within the current fold's training split}.
Sampling proceeds with \texttt{replacement=True}, drawing as many indices per
epoch as the split contains, which makes each represented class equally likely
in expectation while leaving individual batches uneven. This changes class
exposure during training, not the number of distinct recorded expressions.
It is also what makes the class-weighting rule well founded rather than merely
convenient: because the skew correction genuinely happens here, disabling its
loss-side equivalent avoids double-counting rather than removing the
correction (\autoref{eq:focal}).

\textbf{Parameter counts} (\autoref{sec:parameter-counts}). Four components
combine to give any configuration's total: the convolutional backbone
(\texttt{ThreeStreamCNNBackbone}, all three streams) \textbf{14,544}; the
transformer encoder (\texttt{SLSTTTransformer}) \textbf{348,736};
\texttt{RawPatchEmbedding}, the no-CNN fallback, \textbf{4,704}; and the
classifier head \textbf{483}. Their sums give the four distinct totals of
\autoref{tab:parameter-totals}. Neither SimAM nor EVM appears as a row: SimAM
contributes zero parameters by construction (\autoref{eq:simam}), and
magnification is a preprocessing step rather than a network component. Of the
transformer's total, 192 belong to the final \texttt{LayerNorm(96)} applied
after its four encoder layers.


% =============================================================================
\section{Chapter 2 --- Background}
\label{sec:eq-ch2}
% =============================================================================

% -----------------------------------------------------------------------------
\subsection{Eulerian magnification: the small-motion approximation}
\label{eq:evm-taylor}
% -----------------------------------------------------------------------------

\textbf{Appears at} \autoref{sec:eulerian-idea}.

\[
    I(x,t)=f(x+\delta(t))
       \approx f(x)+\delta(t)\frac{\partial f}{\partial x}.
\]

\textbf{Symbols.} $f$ is the image intensity profile; $x$ is spatial
position; $\delta(t)$ is a small time-varying displacement; $I(x,t)$ is the
observed intensity at a \emph{fixed} location $x$ at time $t$.

\textbf{What it does.} It is the first-order Taylor expansion that licenses
the entire Eulerian approach. Eulerian magnification examines intensity
changes at fixed image locations rather than explicitly tracking moving
points, and this approximation is what connects the two: for sufficiently
small displacement, the intensity change observed at a fixed pixel is
proportional to the displacement times the local spatial gradient. Motion can
therefore be manipulated by manipulating intensity.

\textbf{Source.} Bai et al.~\cite{bai2021}, describing the amplitude-based
mechanism. The technique originates with Wu et al., who carry no bibliography
entry under the policy described above and are attributed through the reviewed
paper that reports them.

\textbf{As implemented.} Yes --- but note that the approximation holds only
for small $\delta(t)$. It does not imply that arbitrary intensity changes
represent facial movement, nor that larger amplification always improves
recognition.

% -----------------------------------------------------------------------------
\subsection{Eulerian magnification: the amplified signal}
\label{eq:evm-amplified}
% -----------------------------------------------------------------------------

\textbf{Appears at} \autoref{sec:eulerian-idea}.

\[
    \widetilde I(x,t)\approx
       f(x)+(1+\alpha)\delta(t)\frac{\partial f}{\partial x}.
\]

\textbf{Symbols.} As above, with $\alpha$ the magnification factor and
$\widetilde I$ the reconstructed signal.

\textbf{What it does.} Adding $\alpha$ times the band-passed motion-related
variation back to the original produces a signal that, under the same
first-order approximation, resembles what would have been observed had the
displacement itself been $(1+\alpha)$ times larger. The mechanism is entirely
photometric; nothing is tracked or warped.

\textbf{Source.} Bai et al.~\cite{bai2021}, as above.

\textbf{As implemented.} \textbf{Not in the temporal filter.} The magnitude
relation above is implemented, but the band-pass stage that isolates
$\delta(t)$ differs: Bai et al. describe a Butterworth filter, whereas this
implementation masks Fourier coefficients outside the selected band and
applies the inverse transform (\autoref{sec:temporal-filtering}). The spatial
decomposition uses a Laplacian pyramid with the separable blur kernel
$[1,4,6,4,1]/16$ and factor-two downsampling. A second and larger divergence
concerns \emph{when} magnification is applied: it runs after subsampling, which
makes the realised pass-band clip-dependent (\autoref{eq:ch3-inline}).

% -----------------------------------------------------------------------------
\subsection{Brightness constancy}
\label{eq:brightness}
% -----------------------------------------------------------------------------

\textbf{Appears at} \autoref{sec:optical-flow}.

\[
    I(x,y,t)=I(x+u\Delta t,y+v\Delta t,t+\Delta t).
\]

\textbf{Symbols.} $I$ is image intensity; $(x,y)$ is spatial position; $t$ is
time; $u$ and $v$ are the horizontal and vertical components of local
velocity, for the continuous derivation.

\textbf{What it does.} It states the assumption on which optical flow rests:
a moving point retains its intensity as it moves. Everything estimated as
``motion'' downstream is motion only insofar as this holds.

\textbf{Source.} OFF-ApexNet, Liong et al.~\cite{liong2019a}.

\textbf{As implemented.} Yes, as the assumption underlying the Farneb\"ack
estimator used here --- with two caveats the thesis states rather than
assumes away. In the discrete pipeline $u$ and $v$ denote displacement in
pixels between selected frames rather than physical velocity per second. And
illumination changes or large inter-frame displacement can violate the
approximation outright, which the clip-dependent effective frame rate
(\autoref{eq:ch3-inline}) makes a live concern for the longer clips.

% -----------------------------------------------------------------------------
\subsection{The optical flow constraint}
\label{eq:flow-constraint}
% -----------------------------------------------------------------------------

\textbf{Appears at} \autoref{sec:optical-flow}.

\[
    I_xu+I_yv+I_t=0.
\]

\textbf{Symbols.} $I_x$, $I_y$ and $I_t$ are the partial derivatives of
intensity with respect to $x$, $y$ and $t$; $u$ and $v$ are as above.

\textbf{What it does.} It is the first-order expansion of brightness
constancy, and it is \textbf{one equation in two unknowns}. This
underdetermination is the \textbf{aperture problem}: a local edge constrains
motion normal to the edge, along the intensity gradient, but leaves motion
parallel to the edge entirely undetermined. Every practical estimator must
therefore add something --- a neighbourhood assumption, a local motion model,
or explicit regularisation. Flow fields are consequently never pure
measurements; they carry the assumptions of whichever estimator produced them.

\textbf{Source.} Liong et al.~\cite{liong2019a}.

\textbf{As implemented.} Yes. The constraint is resolved by Farneb\"ack's
polynomial expansion, following the estimator choice reported by Zhao et
al.~\cite{zhaoS2021}, with OpenCV settings: pyramid scale 0.5, three levels,
window size 15, three iterations, polynomial neighbourhood 5, polynomial
smoothing 1.2, flags 0. Flow is computed for the 32 adjacent pairs among the
33 sampled frames.

% -----------------------------------------------------------------------------
\subsection{The symmetric displacement-gradient tensor}
\label{eq:strain-tensor}
% -----------------------------------------------------------------------------

\textbf{Appears at} \autoref{sec:optical-strain}. Its normal and shear
components are also named at \autoref{sec:strain-tensor} in Chapter~3.

\[
    \varepsilon=\tfrac12\left(\nabla\mathbf{u}
                         +(\nabla\mathbf{u})^{\mathsf T}\right),
    \qquad
    \varepsilon_{xx}=\frac{\partial u}{\partial x},\quad
    \varepsilon_{yy}=\frac{\partial v}{\partial y},
\]
\[
    \varepsilon_{xy}=\varepsilon_{yx}
       =\tfrac12\left(\frac{\partial u}{\partial y}
                    +\frac{\partial v}{\partial x}\right).
\]

\textbf{Symbols.} $\mathbf{u}=[u,v]^{\mathsf T}$ is the displacement field;
$\nabla\mathbf{u}$ is its gradient; the superscript $\mathsf{T}$ denotes
transpose. Symmetrising the gradient makes $\varepsilon_{xy}=\varepsilon_{yx}$
by construction.

\textbf{What it does.} It converts displacement into \emph{deformation}.
Flow describes how much things moved; strain describes how neighbouring
displacements \emph{differ}. The normal components describe local stretching or
compression along each axis, and the off-diagonal terms describe shear. The
consequence that motivates using it at all: a spatially uniform translation
has zero displacement gradient and therefore zero strain, so rigid head
movement is suppressed while differential skin motion survives.

\textbf{Source.} The small-strain approximation used for optical strain,
Shreve et al.~\cite{shreve2011}.

\textbf{As implemented.} Yes, with the discretisation as a stated modelling
choice rather than a neutral detail, since it sets the scale at which local
variation is measured. Shreve et al. approximate the spatial derivatives by
central differences with offsets of approximately two to three pixels; Liong
et al.~\cite{liong2014a} use one-pixel offsets, and this implementation
follows the latter via NumPy's \texttt{gradient}. Two limits are declared
rather than argued away: zero strain under uniform translation does not imply
immunity to head \emph{rotation}, projection effects or flow-estimation error,
and differentiation may amplify noise. The reduction of this tensor to a
scalar is where the thesis diverges from the literature --- see
\autoref{eq:strain-implemented}.

% -----------------------------------------------------------------------------
\subsection{Convolution parameter count}
\label{eq:conv-params}
% -----------------------------------------------------------------------------

\textbf{Appears at} \autoref{sec:convolution}.

\[
    N_{\mathrm{params}}=C_{\mathrm{out}}
       (C_{\mathrm{in}}k_Hk_W+1).
\]

\textbf{Symbols.} $C_{\mathrm{in}}$ and $C_{\mathrm{out}}$ are the input and
output channel counts; $k_H$ and $k_W$ are the kernel height and width; the
$+1$ is one bias per output channel.

\textbf{What it does.} It makes weight sharing quantitative. Each output
channel's kernel spans all input channels over its local window, and the same
kernel is applied at every position --- so the parameter count is
\textbf{independent of image dimensions}. That independence is the point of
convolution: no separate detector is learned for each image location. What
does scale with image size is computation and activation memory, since both
grow with the number of positions processed.

\textbf{Source.} Standard; no corpus source and no citation.

\textbf{As implemented.} Yes. The implemented streams use $(1,3,3)$ kernels
with $(0,1,1)$ padding followed by $(1,2,2)$ max-pooling, so the convolutions
preserve temporal length and learn purely spatial filters while pooling halves
height and width. Because $k_T = 1$, \textbf{no temporal convolution is
performed} --- the use of \texttt{Conv3d} does not by itself establish
otherwise. The complete stream can still depend on multiple frames through
BatchNorm during training and through SimAM whenever it is enabled, both of
which use statistics spanning time.

% -----------------------------------------------------------------------------
\subsection{Batch normalisation}
\label{eq:batchnorm}
% -----------------------------------------------------------------------------

\textbf{Appears at} \autoref{sec:normalisation}.

\[
    \widehat x=\frac{x-\mu_c}{\sqrt{\sigma_c^2+\epsilon}},
    \qquad y=\gamma_c\widehat x+\beta_c.
\]

\textbf{Symbols.} $\mu_c$ and $\sigma_c^2$ are the mean and variance for
channel $c$; $\epsilon$ is a small constant preventing division by zero;
$\gamma_c$ and $\beta_c$ are the learned per-channel scale and shift;
$\widehat x$ is the standardised activation and $y$ the output.

\textbf{What it does.} It rescales each channel's activations to a
standardised distribution and then lets the network learn to undo or reshape
that standardisation through $\gamma_c$ and $\beta_c$ --- so normalisation is
imposed but not forced to be retained.

\textbf{Source.} \textbf{None.} Batch normalisation is Ioffe and Szegedy's,
but no paper in this thesis's corpus introduces or examines it, so it is
presented as standard practice and carries no citation. Adding one would
require a new bibliography entry, which the citation policy of this thesis
excludes.

\textbf{As implemented.} Yes. \texttt{BatchNorm3d} computes training
statistics over $(B,T,H,W)$ for each channel, and evaluation uses running
estimates accumulated during training. These activation statistics are a
distinct thing from the per-sample input standardisation performed by the
dataset loader, and the two should not be conflated. Related regularisers in
the same configuration: \texttt{Dropout3d} removes entire feature channels for
a sample rather than masking individual spatial values, and layer
normalisation --- which standardises a token's feature dimension independently
of other samples --- appears inside the transformer and in the classifier
head.

% -----------------------------------------------------------------------------
\subsection{Scaled dot-product attention}
\label{eq:attention}
% -----------------------------------------------------------------------------

\textbf{Appears at} \autoref{sec:self-attention}.

\[
    \operatorname{Attention}(Q,K,V)
      =\operatorname{softmax}\left(\frac{QK^{\mathsf T}}{\sqrt{d_k}}\right)V.
\]

\textbf{Symbols.} $Q$, $K$ and $V$ are the query, key and value matrices,
each a projection of the input tokens; $d_k$ is the key dimension; the
softmax is taken row-wise.

\textbf{What it does.} Queries and keys determine \emph{which} values
contribute to each output, and by how much: the dot product $QK^{\mathsf T}$
scores every token pair, the softmax turns those scores into weights summing
to one, and the result mixes the values accordingly. The $\sqrt{d_k}$ divisor
moderates dot-product magnitude as the key dimension grows, keeping the
softmax out of its saturated regime. Multiple heads learn projections into
different subspaces, whose outputs are concatenated and linearly mixed.

The property that matters most here is negative: without position
information, attention is \textbf{permutation-equivariant}. Permuting the
input tokens permutes the outputs correspondingly, so the operation cannot
identify temporal order on its own. That is what
\autoref{eq:posenc} exists to supply.

\textbf{Source.} Vaswani et al.'s transformer, attributed through Dosovitskiy
et al.~\cite{dosovitskiy2021}, whose Vision Transformer is the foundation this
thesis reaches the architecture through, together with Zhang et al.'s
SLSTT~\cite{zhang2022} (\autoref{sec:transformer-review}). The original
transformer paper carries no bibliography entry under the stated policy.

\textbf{As implemented.} Yes, within \texttt{nn.TransformerEncoder}, but over
a different axis than in the published SLSTT. Here attention operates over
\textbf{32 time steps} of 96-dimensional features, not over spatial patches:
ViT splits an image into patches and relates them spatially, whereas the
sequence here is produced by collapsing both spatial dimensions with
\texttt{AdaptiveAvgPool3d((32,1,1))}. The configuration is $d_{\text{model}} =
96$, $n_{\text{head}} = 8$, four encoder layers, feed-forward dimension 256,
dropout 0.1, pre-norm ordering, GELU activation.

% -----------------------------------------------------------------------------
\subsection{Label smoothing}
\label{eq:label-smoothing}
% -----------------------------------------------------------------------------

\textbf{Appears at} \autoref{sec:label-smoothing}.

\[
    L_{\mathrm{smooth}}
      =-(1-\epsilon)\log p_t
        -\frac{\epsilon}{C}\sum_{c=1}^{C}\log p_c.
\]

\textbf{Symbols.} $\epsilon$ is the smoothing strength; $C$ is the number of
classes; $p_t$ is the predicted probability of the true class and $p_c$ of
class $c$. The target it corresponds to is
$q_c=(1-\epsilon)\mathbf{1}[c=t]+\epsilon/C$, the one-hot vector blended with
uniform mass.

\textbf{What it does.} It removes the incentive to drive the true-class
probability to exactly one. Because a share $\epsilon$ of the target mass sits
on every class, the loss is minimised by a confident but not saturated
prediction. The effect is on \emph{calibration}, not on class balance: it
discourages extreme confidence uniformly and does not preferentially weight a
rare class.

\textbf{Source.} Named, but not examined, by Dosovitskiy et
al.~\cite{dosovitskiy2021}, who list it among the regularisation parameters
they tune. \textbf{No reviewed micro-expression paper uses it, discusses it,
or reports what it does under class skew.}

\textbf{As implemented.} Yes, with $\epsilon = 0.05$, applied \emph{before}
focal modulation --- the focal factor of \autoref{eq:focal} still uses the
unsmoothed true-class probability. Neither the smoothing strength nor the
focal exponent is varied across the architectural comparisons. Because the
corpus supports the choice neither way, it is declared as a choice rather than
defended as a finding, and it appears as such in the divergence ledger.

% -----------------------------------------------------------------------------
\subsection{Precision, recall and per-class F1}
\label{eq:prf}
% -----------------------------------------------------------------------------

\textbf{Appears at} \autoref{sec:confusion-matrix}. The $F1_c$ term also
appears within \autoref{eq:uf1}.

\[
    P_c=\frac{TP_c}{TP_c+FP_c},\qquad
    R_c=\frac{TP_c}{TP_c+FN_c},\qquad
    F1_c=\frac{2TP_c}{2TP_c+FP_c+FN_c}.
\]

\textbf{Symbols.} In a confusion matrix, entry $(i,j)$ counts examples whose
true class is $i$ and predicted class is $j$; the diagonal holds correct
predictions. For class $c$: $TP_c$ counts correct predictions of $c$, $FP_c$
counts incorrect predictions of $c$, and $FN_c$ counts missed examples of $c$.

\textbf{What it does.} Precision penalises false alarms and recall penalises
misses --- each is trivially maximised at the other's expense, which is why
neither is reported alone. F1 is their harmonic mean, so a high value requires
both to be reasonably high; the form above is the algebraic simplification
that avoids computing $P_c$ and $R_c$ separately.

\textbf{Source.} Standard; no corpus source and no citation.

\textbf{As implemented.} Yes, with one convention made explicit because it is
load-bearing under this fold structure: a degenerate denominator, which occurs
whenever a class is absent from a fold, is assigned \textbf{zero} rather than
the class being omitted from the score. That convention is exactly what
produces the ceiling of \autoref{eq:fold-ceiling}, and exactly what pooling
across folds avoids.

% -----------------------------------------------------------------------------
\subsection{Macro-averaged F1}
\label{eq:macro-f1}
% -----------------------------------------------------------------------------

\textbf{Appears at} \autoref{sec:macro-micro}.

\[
    F1_{\mathrm{macro}}=\frac1C\sum_{c=1}^{C}F1_c.
\]

\textbf{Symbols.} $C$ is the number of classes and $F1_c$ is the per-class F1
of \autoref{eq:prf}.

\textbf{What it does.} It gives every class equal weight regardless of how
many examples it has, which is the whole reason for preferring it here.
Contrast micro F1, which first sums $TP$, $FP$ and $FN$ \emph{across classes}:
in single-label multiclass classification each error supplies one false
positive and one false negative, so micro F1 equals accuracy exactly, and can
therefore be dominated by a frequent class. Macro F1 and accuracy describe
different aspects of performance, and reporting both does not remove the need
to inspect per-class errors directly.

\textbf{Source.} Standard. See et al.~\cite{see2019} define the \emph{pooled}
variant, MEGC's UF1, which is the construction this thesis uses
(\autoref{eq:uf1}).

\textbf{As implemented.} Yes, and it is the primary metric --- but only in
its pooled form. Pooling the confusion matrices across folds before computing
any per-class F1 is a genuinely different estimator from computing macro F1
within each fold and averaging the results (\autoref{sec:pooled-vs-fold}). The
distinction is the subject of \autoref{eq:fold-ceiling}, and no run artefact
stores the pooled quantity under any key: it is derived afterwards by the
reporting scripts under \texttt{tools/}.

%
% It must come after \appendix. Including it as a numbered chapter would shift
% the chapter numbering that chapter3_shortened.tex's divergence ledger depends
% on (it cites ten section numbers literally: 3.1.2, 3.1.3, 3.2.6, 3.4.3,
% 3.4.7, 3.5.7, 3.6.6, 3.7.7, 3.8.6, 3.9.7).
%
% CITATION POLICY: this file adds NO new bibliography entries. Works cited only
% inside the reviewed papers (Vaswani, Lin et al., Ioffe & Szegedy, Wu et al.)
% are attributed in prose to the corpus paper that reports them, per the policy
% stated in the header of bib/thesis.bib. Adding an entry would renumber the
% whole bibliography, since tools/bibliography.py numbers by first appearance.
%
% All labels in this file are prefixed eq: and collide with nothing.
% =============================================================================

\chapter{Mathematical Formulations}
\label{ch:formulas}

This appendix collects every display formula that appears in the body of the
thesis, in one place, with a consistent treatment: the expression as it is
written in its own chapter, the section it belongs to, the meaning of each
symbol, the mechanism it implements, the work it comes from, and whether this
thesis implements it as written.

\section*{How to read this appendix}
\addcontentsline{toc}{section}{How to read this appendix}

\textbf{Ordering.} Sequential by chapter, in the order the mathematics is
introduced: Chapter~3 (\autoref{sec:eq-ch3}), Chapter~4
(\autoref{sec:eq-ch4}), then Chapter~2 (\autoref{sec:eq-ch2}), which carries
the bulk of the derivations. Chapters~1, 5 and~6 contain no display
mathematics.

\textbf{Reproduction.} Every expression is reproduced exactly as it is
typeset in its own chapter --- same symbols, same subscripts, same
conventions. Where two chapters state the same relation in different notation,
both forms are given inside the one entry rather than being silently
reconciled.

\textbf{Each formula appears once.} Chapter~2 develops the mechanism and
Chapter~3 reports the literature's version and any divergence from it, so
several relations occur in both. Those are presented once, at the location
listed first in the ordering above, with a cross-reference recording every
section that states them. No entry is duplicated.

\textbf{Attribution.} Citations use only keys already present in
\texttt{bib/thesis.bib}. Several formulas originate in works that the
bibliography deliberately omits, because this thesis's corpus is the set of
reviewed micro-expression papers rather than the full ancestry of every
technique they use. Those are attributed in prose to the reviewed paper that
brings them to this task --- focal loss through Zhao et
al.~\cite{zhaoS2021}, self-attention and the transformer through Dosovitskiy
et al.~\cite{dosovitskiy2021} and Zhang et al.~\cite{zhang2022}, Eulerian
magnification through Bai et al.~\cite{bai2021}. Formulas that are standard
textbook material with no corpus source, such as batch normalisation and the
convolution parameter count, carry no citation and are marked as such.

\textbf{Implementation status.} Each entry closes by stating whether the
thesis implements the formula as written. Three do not, and those are declared
divergences recorded in the divergence ledger of \autoref{tab:divergence-ledger}
rather than defects to be corrected here.

\begin{table}[htbp]
	\centering
	\footnotesize
	\begin{tabular}{@{}l l l l@{}}
		\hline
		\textbf{Entry} & \textbf{Quantity} & \textbf{Attributed to} & \textbf{As written?} \\
		\hline
		\multicolumn{4}{@{}l}{\textit{Chapter 3 --- Literature Review}} \\
		\autoref{eq:uf1}           & Per-class F1 and UF1          & \cite{see2019}                        & yes \\
		\autoref{eq:strain-frobenius} & Four-term strain magnitude & \cite{liong2014a,liong2014b,liong2016} & \textbf{no} \\
		\autoref{eq:strain-implemented} & Three-term strain magnitude & declared divergence                & --- \\
		\autoref{eq:simam}         & SimAM energy and gate         & \cite{yang2021}                       & partly \\
		\autoref{eq:focal}         & Focal loss                    & \cite{zhaoS2021}                      & partly \\
		\multicolumn{4}{@{}l}{\textit{Chapter 4 --- Methodology}} \\
		\autoref{eq:posenc}        & Sinusoidal positional encoding & no corpus source                     & yes \\
		\autoref{eq:fold-ceiling}  & Mean-of-folds macro F1 ceiling & this thesis                          & n/a \\
		\multicolumn{4}{@{}l}{\textit{Chapter 2 --- Background}} \\
		\autoref{eq:evm-taylor}    & Small-motion approximation    & \cite{bai2021}                        & yes \\
		\autoref{eq:evm-amplified} & Amplified signal              & \cite{bai2021}                        & \textbf{no} \\
		\autoref{eq:brightness}    & Brightness constancy          & \cite{liong2019a}                     & yes \\
		\autoref{eq:flow-constraint} & Optical flow constraint     & \cite{liong2019a}                     & yes \\
		\autoref{eq:strain-tensor} & Symmetric displacement gradient & \cite{shreve2011}                   & yes \\
		\autoref{eq:conv-params}   & Convolution parameter count   & standard                              & yes \\
		\autoref{eq:batchnorm}     & Batch normalisation           & standard, no corpus source            & yes \\
		\autoref{eq:attention}     & Scaled dot-product attention  & via \cite{dosovitskiy2021}            & yes \\
		\autoref{eq:label-smoothing} & Label smoothing             & named by \cite{dosovitskiy2021}       & yes \\
		\autoref{eq:prf}           & Precision, recall, F1         & standard                              & yes \\
		\autoref{eq:macro-f1}      & Macro-averaged F1             & standard; \cite{see2019} pooled       & yes \\
		\hline
	\end{tabular}
	\caption[Index of formulations]{Index of the formulations collected in this
	appendix, with the work each is attributed to and whether this thesis
	implements it as written.}
	\label{tab:formula-index}
\end{table}


% =============================================================================
\section{Chapter 3 --- Literature Review}
\label{sec:eq-ch3}
% =============================================================================

% -----------------------------------------------------------------------------
\subsection{Per-class F1 and Unweighted F1}
\label{eq:uf1}
% -----------------------------------------------------------------------------

\textbf{Appears at} \autoref{sec:megc-metrics}.

\[
	F1_c = \frac{2 \cdot TP_c}{2 \cdot TP_c + FP_c + FN_c}, \qquad \text{UF1} = \frac{1}{C}\sum_c F1_c
\]

\textbf{Symbols.} $TP_c$, $FP_c$ and $FN_c$ are the true positives, false
positives and false negatives for class $c$; $C$ is the number of classes,
three here; $F1_c$ is the harmonic mean of precision and recall for class $c$,
written in the algebraically equivalent form that avoids computing either
separately.

\textbf{What it does.} UF1 scores a classifier without letting a frequent
class dominate, by computing a score per class and averaging the scores with
equal weight rather than averaging over examples. The construction that
matters is not visible in the algebra: the counts $TP_c$, $FP_c$ and $FN_c$
are \emph{accumulated across all $k$ folds before any $F1_c$ is computed}.
That pooling is what distinguishes UF1 from the average of per-fold macro F1
scores, which is a different estimator with a different expectation
(\autoref{sec:pooled-vs-fold}, and \autoref{eq:fold-ceiling} for what it costs).

\textbf{Source.} The metric convention of the Second Facial Micro-Expression
Grand Challenge, See et al.~\cite{see2019}, which mandates UF1 and Unweighted
Average Recall in place of plain accuracy on the explicit grounds of class
imbalance.

\textbf{As implemented.} Yes, and it is the thesis's primary metric. Pooled
macro F1 as computed here is identical in construction to MEGC's UF1. The
thesis adopts the challenge's metric definitions while declining its composite
training regime.

% -----------------------------------------------------------------------------
\subsection{Four-term strain magnitude (the published convention)}
\label{eq:strain-frobenius}
% -----------------------------------------------------------------------------

\textbf{Appears at} \autoref{sec:strain-tensor}.

\[
    \varepsilon_{\mathrm{F}}
    =\sqrt{\varepsilon_{xx}^{2}+\varepsilon_{yy}^{2}
        +\varepsilon_{xy}^{2}+\varepsilon_{yx}^{2}}
    =\sqrt{\varepsilon_{xx}^{2}+\varepsilon_{yy}^{2}
        +2\varepsilon_{xy}^{2}}.
\]

\textbf{Symbols.} $\varepsilon_{xx}$ and $\varepsilon_{yy}$ are the normal
components of the symmetric displacement-gradient tensor, describing local
stretching or compression along each axis; $\varepsilon_{xy}=\varepsilon_{yx}$
are its shear components. The tensor itself is defined at
\autoref{eq:strain-tensor}. The second equality follows from the tensor's
symmetry, which makes the two off-diagonal entries equal.

\textbf{What it does.} It reduces the four-component tensor at each pixel to
one non-negative scalar, the Frobenius norm, describing deformation
\emph{intensity} while discarding its sign and direction. Because the tensor is
symmetric, the norm counts shear twice --- once for each off-diagonal entry ---
which is what produces the factor of two.

\textbf{Source.} The Liong strain series~\cite{liong2014a,liong2014b,liong2016}.
Shreve et al.~\cite{shreve2011} establish the tensor; the Liong papers use this
four-component magnitude as the descriptor.

\textbf{As implemented.} \textbf{No.} The implementation computes the
three-term form of \autoref{eq:strain-implemented} instead. This is the
thesis's most consequential divergence and is treated in the next entry.

% -----------------------------------------------------------------------------
\subsection{Three-term strain magnitude (the implemented form)}
\label{eq:strain-implemented}
% -----------------------------------------------------------------------------

\textbf{Appears at} \autoref{sec:strain-divergence}, and again as
$\varepsilon_{\mathrm{mag}}$ at \autoref{sec:optical-strain} in Chapter~2.

\[
    \varepsilon_{\mathrm{os}}
    =\sqrt{\varepsilon_{xx}^{2}+\varepsilon_{yy}^{2}
        +\varepsilon_{xy}^{2}},
    \qquad
    \varepsilon_{xy}=\tfrac{1}{2}
        \left(\frac{\partial u}{\partial y}
             +\frac{\partial v}{\partial x}\right).
\]

Chapter~2 writes the same scalar as

\[
    \varepsilon_{\mathrm{mag}}
       =\sqrt{\varepsilon_{xx}^2+\varepsilon_{yy}^2
                    +\varepsilon_{xy}^2}.
\]

\textbf{Symbols.} As \autoref{eq:strain-frobenius}, with $u$ and $v$ the
horizontal and vertical components of the estimated flow field. In the
discrete pipeline $u$ and $v$ are displacements in pixels between selected
frames, not physical velocities.

\textbf{What it does.} The same reduction to a deformation-intensity scalar,
but counting the shear contribution once rather than twice. The consequence is
a change in relative weighting, not a removable constant: the four-term form
weights pure shear by $\sqrt{2}$ relative to this one. Independent min--max
normalisation removes a uniform scale factor but does not generally remove
this difference, because normal and shear contributions vary across pixels and
across time.

\textbf{Source.} \textbf{None --- this is a declared divergence.} STSTNet
\cite{liong2019b} cannot establish a precedent for the three-term form,
because its printed magnitude equation is typographically defective in two
independent ways. It repeats $\partial u$ in the mixed derivative where the
symmetric tensor calls for $\partial v / \partial x$ alongside $\partial u /
\partial y$, and it places the factor $\tfrac{1}{2}$ \emph{outside} the
squared derivative sum. Correcting the repeated derivative while retaining
that placement gives

\[
    \tfrac{1}{2}
       \left(\frac{\partial u}{\partial y}
            +\frac{\partial v}{\partial x}\right)^2
       =2\varepsilon_{xy}^{2},
\]

which agrees with the \emph{four}-term convention, not the three-term one.
Moving the factor inside the square would instead yield the implemented form,
but that is a further change to the printed expression. The intended
implementation cannot be settled from the published equation alone, so the
thesis declares a departure from the unambiguous four-term convention rather
than claiming to reproduce STSTNet's magnitude.

\textbf{As implemented.} This is what the code computes. Spatial derivatives
come from NumPy's \texttt{gradient} at one-pixel spacing, using central
differences in the array interior; the spacing follows Liong et
al.~\cite{liong2014a} even though the scalar reduction does not. The two
conventions were never compared experimentally here, so their practical
equivalence is not assumed. This is row 4 of the divergence ledger
(\autoref{tab:divergence-ledger}). \textbf{Neither expression should be
``corrected'' --- the difference between them is the finding.}

% -----------------------------------------------------------------------------
\subsection{SimAM: energy and gate}
\label{eq:simam}
% -----------------------------------------------------------------------------

\textbf{Appears at} \autoref{sec:simam-energy}, and at \autoref{sec:simam} in
Chapter~2. Chapter~3 writes it in terms of the intermediate $d_i$ and $s^2$:

\[
    d_i=(x_i-\mu)^2,\qquad
    s^2=\frac{1}{M-1}\sum_{i=1}^{M}d_i,\qquad
    E_i^{-1}=\frac{d_i}{4(s^2+\lambda)}+\frac{1}{2},
\]
\[
    \widetilde{x}_i=x_i\,\operatorname{sigmoid}(E_i^{-1}),
    \qquad \mu=\frac{1}{M}\sum_{i=1}^{M}x_i.
\]

Chapter~2 states the same gate with the variance written as $v$:

\[
    \mu=\frac1M\sum_i x_i,\qquad
    v=\frac1{M-1}\sum_i(x_i-\mu)^2,
\]
\[
    E_i^{-1}=\frac{(x_i-\mu)^2}{4(v+\lambda)}+\frac12,
    \qquad \widetilde x_i=x_i\operatorname{sigmoid}(E_i^{-1}).
\]

\textbf{Symbols.} $x_i$ is one activation; $\mu$ and $s^2$ (equivalently $v$)
are the mean and variance over the statistical context shared by that
channel; $M$ is the number of positions in that context; $\lambda$ is a
positive regularisation constant, fixed at $10^{-4}$ here; $E_i^{-1}$ is the
reciprocal energy, the importance score; $\widetilde{x}_i$ is the refined
activation.

\textbf{What it does.} It implements a distinctiveness prior: an activation
that differs from its channel context is scored as more important, and the
score multiplicatively rescales the feature map rather than being added to it.
The score is the closed-form minimiser of a regularised separation energy, so
the $+\tfrac{1}{2}$ follows from the reciprocal-energy expression rather than
being an arbitrary offset. The sigmoid bounds the multiplicative weight while
preserving the ordering of the scores; it does not produce a hard mask of
inactive regions. No weight or bias is learned anywhere in the module, so it
contributes \textbf{zero learnable parameters} regardless of channel count.

Note what the prior is and is not. It is a statement about activation
statistics, not a detector of facial action units or of minority classes. A
distinctive response may encode useful deformation, or noise, or an artefact.

\textbf{Source.} Yang et al.~\cite{yang2021}.

\textbf{As implemented.} Partly --- with two documented distinctions, both
deliberate.

\begin{enumerate}
	\item \textbf{The variance denominator.} The shared variance printed with
	Yang et al.'s equation~(5) divides by $M$; the reference code in their
	Figure~3 divides by $M-1$. The forms above follow the \emph{code}, as does
	this thesis, which applies Bessel's correction over $n = D \cdot H \cdot W
	- 1$ (\autoref{sec:simam-placement}).
	\item \textbf{The statistical context.} Here $M = THW$ --- whole-clip
	statistics within each channel --- against the original image
	formulation's $M = HW$ per frame. The original module produces a separate
	weight at each channel--height--width position, so its ``three-dimensional
	weights'' describe $C \times H \times W$ and not a temporal volume.
	Extending the context across frames is an additional design decision, not
	part of the source formulation, and is row 10 of the divergence ledger.
	Whole-clip statistics do not guarantee that apex frames receive the
	highest weights.
\end{enumerate}

The value $\lambda = 10^{-4}$ is likewise inherited rather than derived: Yang
et al. sweep $\lambda$ from $10^{-1}$ to $10^{-6}$ and select $10^{-4}$ for
CIFAR, repeating the selection separately for ImageNet and choosing $0.1$. It
is a dataset-selected operating point, not a universal setting justified by
the formula.

% -----------------------------------------------------------------------------
\subsection{Focal loss}
\label{eq:focal}
% -----------------------------------------------------------------------------

\textbf{Appears at} \autoref{sec:imbalance-focal}, and at
\autoref{sec:focal-loss} in Chapter~2.

\[
	\text{FL}(p_t) = -\alpha_t (1 - p_t)^{\gamma}\log p_t
\]

\textbf{Symbols.} $p_t$ is the model's predicted probability for the true
class; $\gamma$ is the focusing exponent, described by Zhao et al. as ``the
balance factor for loss''; $\alpha_t$ is a per-class multiplier, ``the weight
balance factor for samples''. With $\alpha_t = 1$ and $\gamma = 0$ the
expression reduces to ordinary cross-entropy, $-\log p_t$.

\textbf{What it does.} Two independent corrections, which are easy to
conflate and which this thesis keeps separate deliberately. The modulating
factor $(1-p_t)^{\gamma}$ re-weights by \textbf{difficulty}: for $\gamma > 0$
it shrinks the contribution of confident, correct examples relative to hard
ones. It does not inspect class frequency at all, so a difficult example from
any class retains a large contribution --- though under skew the examples it
up-weights are disproportionately from the minority classes. The separate
$\alpha_t$ multiplier re-weights by \textbf{class frequency}. Zhao et al.
treat $\gamma$ as a constant and $\alpha$ as dataset-specific, to be set by
cross-validation.

\textbf{Source.} Zhao et al.~\cite{zhaoS2021}, who bring focal loss into
micro-expression recognition ``to alleviate the inefficient model training
caused by the class imbalance in the MEs Datasets''. They adopt the
formulation developed for one-stage object detection, where the imbalance
being corrected is between foreground and background regions rather than
between emotion classes. Per the bibliography's stated policy, the original
detection paper carries no entry here and the formulation is attributed to the
reviewed paper that reports it.

\textbf{As implemented.} Partly. The difficulty term is used as written, with
$\gamma = 2$ fixed throughout, matching Zhao et al.'s reported exponent --- an
inherited setting, not a measured optimum for this corpus. The implementation
computes $p_t$ from the hard target label using log-softmax, then multiplies a
\emph{label-smoothed} cross-entropy term (\autoref{eq:label-smoothing}) by
$(1-p_t)^{\gamma}$, with the class weight applied afterwards and optionally.
The frequency term $\alpha_t$ is where this thesis departs: frequency
correction is assigned to the sampler instead, and when balanced sampling is
active a runtime guard disables the loss-side class weights, leaving focal
difficulty weighting in place. Conflating the two factors is precisely the
error recorded at \autoref{sec:imbalance-double-correction}.

% -----------------------------------------------------------------------------
\subsection{Inline quantities}
\label{eq:ch3-inline}
% -----------------------------------------------------------------------------

Three quantities in Chapter~3 are not set as display mathematics but carry the
same evidential weight.

\textbf{Flow magnitude, $\sqrt{u^2+v^2}$} (\autoref{sec:brightness-constancy}).
The scalar magnitude of the flow field, described by Zhao et
al.~\cite{zhaoS2021} alongside the component fields $u(x,y)$ and $v(x,y)$.
This thesis retains the two components as separate input channels rather than
reducing them to this magnitude, so the direction of motion is preserved.

\textbf{Effective frame rate, $6600/N$ fps} (\autoref{sec:tim-implications}).
Sampling 33 frames from a clip of $N$ frames recorded at 200~fps yields this
effective rate: \textbf{100~fps} at the median clip length $N=66$, which is
SMIC's recording rate and half what CASME~II was built to provide, and roughly
\textbf{52~fps} for the longest clip at $N=126$, below the original CASME
database's 60~fps~\cite{yan2014}. Eighty-one of the 156 clips are downsampled
by a factor of two or more. No reviewed paper reports this cost.

\textbf{True inter-frame interval, $N/(200 \times 32)$ seconds}
(\autoref{sec:tim-implications}). Magnification runs on the 33 sampled frames
at the nominal 200~fps, but those frames span the clip's full duration, so the
real interval between them is this rather than $1/200$. A band specified as
5--25~Hz therefore corresponds to roughly 2.5--12.5~Hz physically at the
median, and the discrepancy grows with clip length. The realised pass-band is
consequently clip-dependent --- neither the 5--25~Hz derived from the duration
rule of Bai et al.~\cite{bai2021} nor constant across the corpus. This is a
declared divergence; magnifying before subsampling is declared as further
work.


% =============================================================================
\section{Chapter 4 --- Methodology}
\label{sec:eq-ch4}
% =============================================================================

% -----------------------------------------------------------------------------
\subsection{Sinusoidal positional encoding}
\label{eq:posenc}
% -----------------------------------------------------------------------------

\textbf{Appears at} \autoref{sec:transformer-config}.

\[
	\mathrm{PE}(p, 2i) = \sin\!\left(\frac{p}{10000^{2i/96}}\right), \qquad \mathrm{PE}(p, 2i+1) = \cos\!\left(\frac{p}{10000^{2i/96}}\right).
\]

\textbf{Symbols.} $p$ is the position in the sequence, one of the 32 time
steps; $i$ indexes feature pairs within the embedding, so $2i$ and $2i+1$ are
the even and odd feature positions; 96 is $d_{\text{model}}$, the embedding
width, appearing here as a literal rather than a symbol because the
configuration fixes it.

\textbf{What it does.} It supplies the order information that self-attention
cannot recover on its own. Without it, attention is permutation-equivariant:
permuting the input tokens permutes the outputs correspondingly, so a shuffled
clip is indistinguishable from an ordered one (\autoref{sec:self-attention}).
Each feature pair oscillates at a different frequency, geometrically spaced by
the $10000^{2i/96}$ denominator, so the vector at each position is distinct
and positions at a fixed offset stand in a consistent linear relation.

\textbf{Source.} No corpus source. The encoding is Vaswani et al.'s, from the
original transformer, which the bibliography deliberately omits under the
policy described at the head of this appendix; Chapter~4 states it as the
standard form without attribution. The thesis reaches the architecture through
Zhang et al.'s SLSTT~\cite{zhang2022} and its Vision Transformer
foundation~\cite{dosovitskiy2021}, as traced at
\autoref{sec:transformer-review}. Note that ViT itself uses \emph{learned}
positional embeddings, not these fixed sinusoids.

\textbf{As implemented.} Yes, as written. It is a fixed buffer, not a learned
parameter, and contributes nothing to any parameter count. Two details of the
surrounding configuration are easily missed: the positional-encoding stage
applies its own \texttt{Dropout(p=0.1)} to the sequence \emph{after} the
encoding is added, which is a separate application from the encoder layers'
dropout at the same rate; and although the sinusoids could be evaluated at
arbitrary positions, the implementation uses a finite precomputed buffer, so
generalisation to unseen sequence lengths is neither exercised nor guaranteed.

% -----------------------------------------------------------------------------
\subsection{The mean-of-folds macro F1 ceiling}
\label{eq:fold-ceiling}
% -----------------------------------------------------------------------------

\textbf{Appears at} \autoref{sec:mean-of-folds-rejected}.

\[
	\frac{10 \cdot \tfrac13 + 8 \cdot \tfrac23 + 7 \cdot 1}{25} = \frac{3.333\ldots + 5.333\ldots + 7}{25} \approx 0.627.
\]

\textbf{Symbols.} The three numerator terms are the 25 leave-one-subject-out
folds grouped by how many of the three classes appear in each fold's held-out
ground truth: \textbf{10} folds contain one class, \textbf{8} contain two,
\textbf{7} contain all three. The fractions $\tfrac13$, $\tfrac23$ and $1$ are
the maximum macro F1 each group can attain. The denominator is the fold count.

\textbf{What it does.} It computes an upper bound, fixed by protocol
structure before any model is trained. A macro F1 computed \emph{within} a
fold is the unweighted mean of that fold's three per-class F1 scores, and a
class absent from that fold's ground truth contributes zero no matter what the
model predicts. A single-class fold is therefore capped at $1/3$ and a
two-class fold at $2/3$. Averaging across folds averages seventeen sub-unity
ceilings into every score, bounding the mean-of-folds macro F1 at
approximately \textbf{0.627} --- a property of which subjects were held out,
identical for all twelve configurations.

Pooling avoids the ceiling entirely: accumulating $TP$, $FP$ and $FN$ across
all 25 folds and all 156 clips before computing any per-class F1 means a class
locally absent from one fold no longer forces a zero into its contribution.
That is the construction of \autoref{eq:uf1}.

\textbf{Source.} This thesis's own arithmetic over the fold composition of
\autoref{fig:fold-composition}. It is not a published result.

\textbf{As implemented.} Not applicable --- this is an argument for rejecting
a metric, not a component. Its consequence is directly implemented: pooled
macro F1 is the primary reported metric, and the mean-of-folds quantity stored
under \texttt{macro\_f1} in the run artefacts is not. A metric with a hard
ceiling below 0.7 cannot serve as a headline result. The same fold composition
biases mean-of-folds \emph{accuracy} in the same direction without imposing a
hard ceiling; for the full configuration it exceeds pooled accuracy by
\textbf{6.3 points}.

% -----------------------------------------------------------------------------
\subsection{Inline quantities}
\label{eq:ch4-inline}
% -----------------------------------------------------------------------------

\textbf{Sampler weight, $1/n_c$} (\autoref{sec:sampler-config}). Every
training example is assigned this weight by the
\texttt{WeightedRandomSampler}, where $n_c$ is the number of training examples
of that example's class \emph{within the current fold's training split}.
Sampling proceeds with \texttt{replacement=True}, drawing as many indices per
epoch as the split contains, which makes each represented class equally likely
in expectation while leaving individual batches uneven. This changes class
exposure during training, not the number of distinct recorded expressions.
It is also what makes the class-weighting rule well founded rather than merely
convenient: because the skew correction genuinely happens here, disabling its
loss-side equivalent avoids double-counting rather than removing the
correction (\autoref{eq:focal}).

\textbf{Parameter counts} (\autoref{sec:parameter-counts}). Four components
combine to give any configuration's total: the convolutional backbone
(\texttt{ThreeStreamCNNBackbone}, all three streams) \textbf{14,544}; the
transformer encoder (\texttt{SLSTTTransformer}) \textbf{348,736};
\texttt{RawPatchEmbedding}, the no-CNN fallback, \textbf{4,704}; and the
classifier head \textbf{483}. Their sums give the four distinct totals of
\autoref{tab:parameter-totals}. Neither SimAM nor EVM appears as a row: SimAM
contributes zero parameters by construction (\autoref{eq:simam}), and
magnification is a preprocessing step rather than a network component. Of the
transformer's total, 192 belong to the final \texttt{LayerNorm(96)} applied
after its four encoder layers.


% =============================================================================
\section{Chapter 2 --- Background}
\label{sec:eq-ch2}
% =============================================================================

% -----------------------------------------------------------------------------
\subsection{Eulerian magnification: the small-motion approximation}
\label{eq:evm-taylor}
% -----------------------------------------------------------------------------

\textbf{Appears at} \autoref{sec:eulerian-idea}.

\[
    I(x,t)=f(x+\delta(t))
       \approx f(x)+\delta(t)\frac{\partial f}{\partial x}.
\]

\textbf{Symbols.} $f$ is the image intensity profile; $x$ is spatial
position; $\delta(t)$ is a small time-varying displacement; $I(x,t)$ is the
observed intensity at a \emph{fixed} location $x$ at time $t$.

\textbf{What it does.} It is the first-order Taylor expansion that licenses
the entire Eulerian approach. Eulerian magnification examines intensity
changes at fixed image locations rather than explicitly tracking moving
points, and this approximation is what connects the two: for sufficiently
small displacement, the intensity change observed at a fixed pixel is
proportional to the displacement times the local spatial gradient. Motion can
therefore be manipulated by manipulating intensity.

\textbf{Source.} Bai et al.~\cite{bai2021}, describing the amplitude-based
mechanism. The technique originates with Wu et al., who carry no bibliography
entry under the policy described above and are attributed through the reviewed
paper that reports them.

\textbf{As implemented.} Yes --- but note that the approximation holds only
for small $\delta(t)$. It does not imply that arbitrary intensity changes
represent facial movement, nor that larger amplification always improves
recognition.

% -----------------------------------------------------------------------------
\subsection{Eulerian magnification: the amplified signal}
\label{eq:evm-amplified}
% -----------------------------------------------------------------------------

\textbf{Appears at} \autoref{sec:eulerian-idea}.

\[
    \widetilde I(x,t)\approx
       f(x)+(1+\alpha)\delta(t)\frac{\partial f}{\partial x}.
\]

\textbf{Symbols.} As above, with $\alpha$ the magnification factor and
$\widetilde I$ the reconstructed signal.

\textbf{What it does.} Adding $\alpha$ times the band-passed motion-related
variation back to the original produces a signal that, under the same
first-order approximation, resembles what would have been observed had the
displacement itself been $(1+\alpha)$ times larger. The mechanism is entirely
photometric; nothing is tracked or warped.

\textbf{Source.} Bai et al.~\cite{bai2021}, as above.

\textbf{As implemented.} \textbf{Not in the temporal filter.} The magnitude
relation above is implemented, but the band-pass stage that isolates
$\delta(t)$ differs: Bai et al. describe a Butterworth filter, whereas this
implementation masks Fourier coefficients outside the selected band and
applies the inverse transform (\autoref{sec:temporal-filtering}). The spatial
decomposition uses a Laplacian pyramid with the separable blur kernel
$[1,4,6,4,1]/16$ and factor-two downsampling. A second and larger divergence
concerns \emph{when} magnification is applied: it runs after subsampling, which
makes the realised pass-band clip-dependent (\autoref{eq:ch3-inline}).

% -----------------------------------------------------------------------------
\subsection{Brightness constancy}
\label{eq:brightness}
% -----------------------------------------------------------------------------

\textbf{Appears at} \autoref{sec:optical-flow}.

\[
    I(x,y,t)=I(x+u\Delta t,y+v\Delta t,t+\Delta t).
\]

\textbf{Symbols.} $I$ is image intensity; $(x,y)$ is spatial position; $t$ is
time; $u$ and $v$ are the horizontal and vertical components of local
velocity, for the continuous derivation.

\textbf{What it does.} It states the assumption on which optical flow rests:
a moving point retains its intensity as it moves. Everything estimated as
``motion'' downstream is motion only insofar as this holds.

\textbf{Source.} OFF-ApexNet, Liong et al.~\cite{liong2019a}.

\textbf{As implemented.} Yes, as the assumption underlying the Farneb\"ack
estimator used here --- with two caveats the thesis states rather than
assumes away. In the discrete pipeline $u$ and $v$ denote displacement in
pixels between selected frames rather than physical velocity per second. And
illumination changes or large inter-frame displacement can violate the
approximation outright, which the clip-dependent effective frame rate
(\autoref{eq:ch3-inline}) makes a live concern for the longer clips.

% -----------------------------------------------------------------------------
\subsection{The optical flow constraint}
\label{eq:flow-constraint}
% -----------------------------------------------------------------------------

\textbf{Appears at} \autoref{sec:optical-flow}.

\[
    I_xu+I_yv+I_t=0.
\]

\textbf{Symbols.} $I_x$, $I_y$ and $I_t$ are the partial derivatives of
intensity with respect to $x$, $y$ and $t$; $u$ and $v$ are as above.

\textbf{What it does.} It is the first-order expansion of brightness
constancy, and it is \textbf{one equation in two unknowns}. This
underdetermination is the \textbf{aperture problem}: a local edge constrains
motion normal to the edge, along the intensity gradient, but leaves motion
parallel to the edge entirely undetermined. Every practical estimator must
therefore add something --- a neighbourhood assumption, a local motion model,
or explicit regularisation. Flow fields are consequently never pure
measurements; they carry the assumptions of whichever estimator produced them.

\textbf{Source.} Liong et al.~\cite{liong2019a}.

\textbf{As implemented.} Yes. The constraint is resolved by Farneb\"ack's
polynomial expansion, following the estimator choice reported by Zhao et
al.~\cite{zhaoS2021}, with OpenCV settings: pyramid scale 0.5, three levels,
window size 15, three iterations, polynomial neighbourhood 5, polynomial
smoothing 1.2, flags 0. Flow is computed for the 32 adjacent pairs among the
33 sampled frames.

% -----------------------------------------------------------------------------
\subsection{The symmetric displacement-gradient tensor}
\label{eq:strain-tensor}
% -----------------------------------------------------------------------------

\textbf{Appears at} \autoref{sec:optical-strain}. Its normal and shear
components are also named at \autoref{sec:strain-tensor} in Chapter~3.

\[
    \varepsilon=\tfrac12\left(\nabla\mathbf{u}
                         +(\nabla\mathbf{u})^{\mathsf T}\right),
    \qquad
    \varepsilon_{xx}=\frac{\partial u}{\partial x},\quad
    \varepsilon_{yy}=\frac{\partial v}{\partial y},
\]
\[
    \varepsilon_{xy}=\varepsilon_{yx}
       =\tfrac12\left(\frac{\partial u}{\partial y}
                    +\frac{\partial v}{\partial x}\right).
\]

\textbf{Symbols.} $\mathbf{u}=[u,v]^{\mathsf T}$ is the displacement field;
$\nabla\mathbf{u}$ is its gradient; the superscript $\mathsf{T}$ denotes
transpose. Symmetrising the gradient makes $\varepsilon_{xy}=\varepsilon_{yx}$
by construction.

\textbf{What it does.} It converts displacement into \emph{deformation}.
Flow describes how much things moved; strain describes how neighbouring
displacements \emph{differ}. The normal components describe local stretching or
compression along each axis, and the off-diagonal terms describe shear. The
consequence that motivates using it at all: a spatially uniform translation
has zero displacement gradient and therefore zero strain, so rigid head
movement is suppressed while differential skin motion survives.

\textbf{Source.} The small-strain approximation used for optical strain,
Shreve et al.~\cite{shreve2011}.

\textbf{As implemented.} Yes, with the discretisation as a stated modelling
choice rather than a neutral detail, since it sets the scale at which local
variation is measured. Shreve et al. approximate the spatial derivatives by
central differences with offsets of approximately two to three pixels; Liong
et al.~\cite{liong2014a} use one-pixel offsets, and this implementation
follows the latter via NumPy's \texttt{gradient}. Two limits are declared
rather than argued away: zero strain under uniform translation does not imply
immunity to head \emph{rotation}, projection effects or flow-estimation error,
and differentiation may amplify noise. The reduction of this tensor to a
scalar is where the thesis diverges from the literature --- see
\autoref{eq:strain-implemented}.

% -----------------------------------------------------------------------------
\subsection{Convolution parameter count}
\label{eq:conv-params}
% -----------------------------------------------------------------------------

\textbf{Appears at} \autoref{sec:convolution}.

\[
    N_{\mathrm{params}}=C_{\mathrm{out}}
       (C_{\mathrm{in}}k_Hk_W+1).
\]

\textbf{Symbols.} $C_{\mathrm{in}}$ and $C_{\mathrm{out}}$ are the input and
output channel counts; $k_H$ and $k_W$ are the kernel height and width; the
$+1$ is one bias per output channel.

\textbf{What it does.} It makes weight sharing quantitative. Each output
channel's kernel spans all input channels over its local window, and the same
kernel is applied at every position --- so the parameter count is
\textbf{independent of image dimensions}. That independence is the point of
convolution: no separate detector is learned for each image location. What
does scale with image size is computation and activation memory, since both
grow with the number of positions processed.

\textbf{Source.} Standard; no corpus source and no citation.

\textbf{As implemented.} Yes. The implemented streams use $(1,3,3)$ kernels
with $(0,1,1)$ padding followed by $(1,2,2)$ max-pooling, so the convolutions
preserve temporal length and learn purely spatial filters while pooling halves
height and width. Because $k_T = 1$, \textbf{no temporal convolution is
performed} --- the use of \texttt{Conv3d} does not by itself establish
otherwise. The complete stream can still depend on multiple frames through
BatchNorm during training and through SimAM whenever it is enabled, both of
which use statistics spanning time.

% -----------------------------------------------------------------------------
\subsection{Batch normalisation}
\label{eq:batchnorm}
% -----------------------------------------------------------------------------

\textbf{Appears at} \autoref{sec:normalisation}.

\[
    \widehat x=\frac{x-\mu_c}{\sqrt{\sigma_c^2+\epsilon}},
    \qquad y=\gamma_c\widehat x+\beta_c.
\]

\textbf{Symbols.} $\mu_c$ and $\sigma_c^2$ are the mean and variance for
channel $c$; $\epsilon$ is a small constant preventing division by zero;
$\gamma_c$ and $\beta_c$ are the learned per-channel scale and shift;
$\widehat x$ is the standardised activation and $y$ the output.

\textbf{What it does.} It rescales each channel's activations to a
standardised distribution and then lets the network learn to undo or reshape
that standardisation through $\gamma_c$ and $\beta_c$ --- so normalisation is
imposed but not forced to be retained.

\textbf{Source.} \textbf{None.} Batch normalisation is Ioffe and Szegedy's,
but no paper in this thesis's corpus introduces or examines it, so it is
presented as standard practice and carries no citation. Adding one would
require a new bibliography entry, which the citation policy of this thesis
excludes.

\textbf{As implemented.} Yes. \texttt{BatchNorm3d} computes training
statistics over $(B,T,H,W)$ for each channel, and evaluation uses running
estimates accumulated during training. These activation statistics are a
distinct thing from the per-sample input standardisation performed by the
dataset loader, and the two should not be conflated. Related regularisers in
the same configuration: \texttt{Dropout3d} removes entire feature channels for
a sample rather than masking individual spatial values, and layer
normalisation --- which standardises a token's feature dimension independently
of other samples --- appears inside the transformer and in the classifier
head.

% -----------------------------------------------------------------------------
\subsection{Scaled dot-product attention}
\label{eq:attention}
% -----------------------------------------------------------------------------

\textbf{Appears at} \autoref{sec:self-attention}.

\[
    \operatorname{Attention}(Q,K,V)
      =\operatorname{softmax}\left(\frac{QK^{\mathsf T}}{\sqrt{d_k}}\right)V.
\]

\textbf{Symbols.} $Q$, $K$ and $V$ are the query, key and value matrices,
each a projection of the input tokens; $d_k$ is the key dimension; the
softmax is taken row-wise.

\textbf{What it does.} Queries and keys determine \emph{which} values
contribute to each output, and by how much: the dot product $QK^{\mathsf T}$
scores every token pair, the softmax turns those scores into weights summing
to one, and the result mixes the values accordingly. The $\sqrt{d_k}$ divisor
moderates dot-product magnitude as the key dimension grows, keeping the
softmax out of its saturated regime. Multiple heads learn projections into
different subspaces, whose outputs are concatenated and linearly mixed.

The property that matters most here is negative: without position
information, attention is \textbf{permutation-equivariant}. Permuting the
input tokens permutes the outputs correspondingly, so the operation cannot
identify temporal order on its own. That is what
\autoref{eq:posenc} exists to supply.

\textbf{Source.} Vaswani et al.'s transformer, attributed through Dosovitskiy
et al.~\cite{dosovitskiy2021}, whose Vision Transformer is the foundation this
thesis reaches the architecture through, together with Zhang et al.'s
SLSTT~\cite{zhang2022} (\autoref{sec:transformer-review}). The original
transformer paper carries no bibliography entry under the stated policy.

\textbf{As implemented.} Yes, within \texttt{nn.TransformerEncoder}, but over
a different axis than in the published SLSTT. Here attention operates over
\textbf{32 time steps} of 96-dimensional features, not over spatial patches:
ViT splits an image into patches and relates them spatially, whereas the
sequence here is produced by collapsing both spatial dimensions with
\texttt{AdaptiveAvgPool3d((32,1,1))}. The configuration is $d_{\text{model}} =
96$, $n_{\text{head}} = 8$, four encoder layers, feed-forward dimension 256,
dropout 0.1, pre-norm ordering, GELU activation.

% -----------------------------------------------------------------------------
\subsection{Label smoothing}
\label{eq:label-smoothing}
% -----------------------------------------------------------------------------

\textbf{Appears at} \autoref{sec:label-smoothing}.

\[
    L_{\mathrm{smooth}}
      =-(1-\epsilon)\log p_t
        -\frac{\epsilon}{C}\sum_{c=1}^{C}\log p_c.
\]

\textbf{Symbols.} $\epsilon$ is the smoothing strength; $C$ is the number of
classes; $p_t$ is the predicted probability of the true class and $p_c$ of
class $c$. The target it corresponds to is
$q_c=(1-\epsilon)\mathbf{1}[c=t]+\epsilon/C$, the one-hot vector blended with
uniform mass.

\textbf{What it does.} It removes the incentive to drive the true-class
probability to exactly one. Because a share $\epsilon$ of the target mass sits
on every class, the loss is minimised by a confident but not saturated
prediction. The effect is on \emph{calibration}, not on class balance: it
discourages extreme confidence uniformly and does not preferentially weight a
rare class.

\textbf{Source.} Named, but not examined, by Dosovitskiy et
al.~\cite{dosovitskiy2021}, who list it among the regularisation parameters
they tune. \textbf{No reviewed micro-expression paper uses it, discusses it,
or reports what it does under class skew.}

\textbf{As implemented.} Yes, with $\epsilon = 0.05$, applied \emph{before}
focal modulation --- the focal factor of \autoref{eq:focal} still uses the
unsmoothed true-class probability. Neither the smoothing strength nor the
focal exponent is varied across the architectural comparisons. Because the
corpus supports the choice neither way, it is declared as a choice rather than
defended as a finding, and it appears as such in the divergence ledger.

% -----------------------------------------------------------------------------
\subsection{Precision, recall and per-class F1}
\label{eq:prf}
% -----------------------------------------------------------------------------

\textbf{Appears at} \autoref{sec:confusion-matrix}. The $F1_c$ term also
appears within \autoref{eq:uf1}.

\[
    P_c=\frac{TP_c}{TP_c+FP_c},\qquad
    R_c=\frac{TP_c}{TP_c+FN_c},\qquad
    F1_c=\frac{2TP_c}{2TP_c+FP_c+FN_c}.
\]

\textbf{Symbols.} In a confusion matrix, entry $(i,j)$ counts examples whose
true class is $i$ and predicted class is $j$; the diagonal holds correct
predictions. For class $c$: $TP_c$ counts correct predictions of $c$, $FP_c$
counts incorrect predictions of $c$, and $FN_c$ counts missed examples of $c$.

\textbf{What it does.} Precision penalises false alarms and recall penalises
misses --- each is trivially maximised at the other's expense, which is why
neither is reported alone. F1 is their harmonic mean, so a high value requires
both to be reasonably high; the form above is the algebraic simplification
that avoids computing $P_c$ and $R_c$ separately.

\textbf{Source.} Standard; no corpus source and no citation.

\textbf{As implemented.} Yes, with one convention made explicit because it is
load-bearing under this fold structure: a degenerate denominator, which occurs
whenever a class is absent from a fold, is assigned \textbf{zero} rather than
the class being omitted from the score. That convention is exactly what
produces the ceiling of \autoref{eq:fold-ceiling}, and exactly what pooling
across folds avoids.

% -----------------------------------------------------------------------------
\subsection{Macro-averaged F1}
\label{eq:macro-f1}
% -----------------------------------------------------------------------------

\textbf{Appears at} \autoref{sec:macro-micro}.

\[
    F1_{\mathrm{macro}}=\frac1C\sum_{c=1}^{C}F1_c.
\]

\textbf{Symbols.} $C$ is the number of classes and $F1_c$ is the per-class F1
of \autoref{eq:prf}.

\textbf{What it does.} It gives every class equal weight regardless of how
many examples it has, which is the whole reason for preferring it here.
Contrast micro F1, which first sums $TP$, $FP$ and $FN$ \emph{across classes}:
in single-label multiclass classification each error supplies one false
positive and one false negative, so micro F1 equals accuracy exactly, and can
therefore be dominated by a frequent class. Macro F1 and accuracy describe
different aspects of performance, and reporting both does not remove the need
to inspect per-class errors directly.

\textbf{Source.} Standard. See et al.~\cite{see2019} define the \emph{pooled}
variant, MEGC's UF1, which is the construction this thesis uses
(\autoref{eq:uf1}).

\textbf{As implemented.} Yes, and it is the primary metric --- but only in
its pooled form. Pooling the confusion matrices across folds before computing
any per-class F1 is a genuinely different estimator from computing macro F1
within each fold and averaging the results (\autoref{sec:pooled-vs-fold}). The
distinction is the subject of \autoref{eq:fold-ceiling}, and no run artefact
stores the pooled quantity under any key: it is derived afterwards by the
reporting scripts under \texttt{tools/}.

%
% It must come after \appendix. Including it as a numbered chapter would shift
% the chapter numbering that chapter3_shortened.tex's divergence ledger depends
% on (it cites ten section numbers literally: 3.1.2, 3.1.3, 3.2.6, 3.4.3,
% 3.4.7, 3.5.7, 3.6.6, 3.7.7, 3.8.6, 3.9.7).
%
% CITATION POLICY: this file adds NO new bibliography entries. Works cited only
% inside the reviewed papers (Vaswani, Lin et al., Ioffe & Szegedy, Wu et al.)
% are attributed in prose to the corpus paper that reports them, per the policy
% stated in the header of bib/thesis.bib. Adding an entry would renumber the
% whole bibliography, since tools/bibliography.py numbers by first appearance.
%
% All labels in this file are prefixed eq: and collide with nothing.
% =============================================================================

\chapter{Mathematical Formulations}
\label{ch:formulas}

This appendix collects every display formula that appears in the body of the
thesis, in one place, with a consistent treatment: the expression as it is
written in its own chapter, the section it belongs to, the meaning of each
symbol, the mechanism it implements, the work it comes from, and whether this
thesis implements it as written.

\section*{How to read this appendix}
\addcontentsline{toc}{section}{How to read this appendix}

\textbf{Ordering.} Sequential by chapter, in the order the mathematics is
introduced: Chapter~3 (\autoref{sec:eq-ch3}), Chapter~4
(\autoref{sec:eq-ch4}), then Chapter~2 (\autoref{sec:eq-ch2}), which carries
the bulk of the derivations. Chapters~1, 5 and~6 contain no display
mathematics.

\textbf{Reproduction.} Every expression is reproduced exactly as it is
typeset in its own chapter --- same symbols, same subscripts, same
conventions. Where two chapters state the same relation in different notation,
both forms are given inside the one entry rather than being silently
reconciled.

\textbf{Each formula appears once.} Chapter~2 develops the mechanism and
Chapter~3 reports the literature's version and any divergence from it, so
several relations occur in both. Those are presented once, at the location
listed first in the ordering above, with a cross-reference recording every
section that states them. No entry is duplicated.

\textbf{Attribution.} Citations use only keys already present in
\texttt{bib/thesis.bib}. Several formulas originate in works that the
bibliography deliberately omits, because this thesis's corpus is the set of
reviewed micro-expression papers rather than the full ancestry of every
technique they use. Those are attributed in prose to the reviewed paper that
brings them to this task --- focal loss through Zhao et
al.~\cite{zhaoS2021}, self-attention and the transformer through Dosovitskiy
et al.~\cite{dosovitskiy2021} and Zhang et al.~\cite{zhang2022}, Eulerian
magnification through Bai et al.~\cite{bai2021}. Formulas that are standard
textbook material with no corpus source, such as batch normalisation and the
convolution parameter count, carry no citation and are marked as such.

\textbf{Implementation status.} Each entry closes by stating whether the
thesis implements the formula as written. Three do not, and those are declared
divergences recorded in the divergence ledger of \autoref{tab:divergence-ledger}
rather than defects to be corrected here.

\begin{table}[htbp]
	\centering
	\footnotesize
	\begin{tabular}{@{}l l l l@{}}
		\hline
		\textbf{Entry} & \textbf{Quantity} & \textbf{Attributed to} & \textbf{As written?} \\
		\hline
		\multicolumn{4}{@{}l}{\textit{Chapter 3 --- Literature Review}} \\
		\autoref{eq:uf1}           & Per-class F1 and UF1          & \cite{see2019}                        & yes \\
		\autoref{eq:strain-frobenius} & Four-term strain magnitude & \cite{liong2014a,liong2014b,liong2016} & \textbf{no} \\
		\autoref{eq:strain-implemented} & Three-term strain magnitude & declared divergence                & --- \\
		\autoref{eq:simam}         & SimAM energy and gate         & \cite{yang2021}                       & partly \\
		\autoref{eq:focal}         & Focal loss                    & \cite{zhaoS2021}                      & partly \\
		\multicolumn{4}{@{}l}{\textit{Chapter 4 --- Methodology}} \\
		\autoref{eq:posenc}        & Sinusoidal positional encoding & no corpus source                     & yes \\
		\autoref{eq:fold-ceiling}  & Mean-of-folds macro F1 ceiling & this thesis                          & n/a \\
		\multicolumn{4}{@{}l}{\textit{Chapter 2 --- Background}} \\
		\autoref{eq:evm-taylor}    & Small-motion approximation    & \cite{bai2021}                        & yes \\
		\autoref{eq:evm-amplified} & Amplified signal              & \cite{bai2021}                        & \textbf{no} \\
		\autoref{eq:brightness}    & Brightness constancy          & \cite{liong2019a}                     & yes \\
		\autoref{eq:flow-constraint} & Optical flow constraint     & \cite{liong2019a}                     & yes \\
		\autoref{eq:strain-tensor} & Symmetric displacement gradient & \cite{shreve2011}                   & yes \\
		\autoref{eq:conv-params}   & Convolution parameter count   & standard                              & yes \\
		\autoref{eq:batchnorm}     & Batch normalisation           & standard, no corpus source            & yes \\
		\autoref{eq:attention}     & Scaled dot-product attention  & via \cite{dosovitskiy2021}            & yes \\
		\autoref{eq:label-smoothing} & Label smoothing             & named by \cite{dosovitskiy2021}       & yes \\
		\autoref{eq:prf}           & Precision, recall, F1         & standard                              & yes \\
		\autoref{eq:macro-f1}      & Macro-averaged F1             & standard; \cite{see2019} pooled       & yes \\
		\hline
	\end{tabular}
	\caption[Index of formulations]{Index of the formulations collected in this
	appendix, with the work each is attributed to and whether this thesis
	implements it as written.}
	\label{tab:formula-index}
\end{table}


% =============================================================================
\section{Chapter 3 --- Literature Review}
\label{sec:eq-ch3}
% =============================================================================

% -----------------------------------------------------------------------------
\subsection{Per-class F1 and Unweighted F1}
\label{eq:uf1}
% -----------------------------------------------------------------------------

\textbf{Appears at} \autoref{sec:megc-metrics}.

\[
	F1_c = \frac{2 \cdot TP_c}{2 \cdot TP_c + FP_c + FN_c}, \qquad \text{UF1} = \frac{1}{C}\sum_c F1_c
\]

\textbf{Symbols.} $TP_c$, $FP_c$ and $FN_c$ are the true positives, false
positives and false negatives for class $c$; $C$ is the number of classes,
three here; $F1_c$ is the harmonic mean of precision and recall for class $c$,
written in the algebraically equivalent form that avoids computing either
separately.

\textbf{What it does.} UF1 scores a classifier without letting a frequent
class dominate, by computing a score per class and averaging the scores with
equal weight rather than averaging over examples. The construction that
matters is not visible in the algebra: the counts $TP_c$, $FP_c$ and $FN_c$
are \emph{accumulated across all $k$ folds before any $F1_c$ is computed}.
That pooling is what distinguishes UF1 from the average of per-fold macro F1
scores, which is a different estimator with a different expectation
(\autoref{sec:pooled-vs-fold}, and \autoref{eq:fold-ceiling} for what it costs).

\textbf{Source.} The metric convention of the Second Facial Micro-Expression
Grand Challenge, See et al.~\cite{see2019}, which mandates UF1 and Unweighted
Average Recall in place of plain accuracy on the explicit grounds of class
imbalance.

\textbf{As implemented.} Yes, and it is the thesis's primary metric. Pooled
macro F1 as computed here is identical in construction to MEGC's UF1. The
thesis adopts the challenge's metric definitions while declining its composite
training regime.

% -----------------------------------------------------------------------------
\subsection{Four-term strain magnitude (the published convention)}
\label{eq:strain-frobenius}
% -----------------------------------------------------------------------------

\textbf{Appears at} \autoref{sec:strain-tensor}.

\[
    \varepsilon_{\mathrm{F}}
    =\sqrt{\varepsilon_{xx}^{2}+\varepsilon_{yy}^{2}
        +\varepsilon_{xy}^{2}+\varepsilon_{yx}^{2}}
    =\sqrt{\varepsilon_{xx}^{2}+\varepsilon_{yy}^{2}
        +2\varepsilon_{xy}^{2}}.
\]

\textbf{Symbols.} $\varepsilon_{xx}$ and $\varepsilon_{yy}$ are the normal
components of the symmetric displacement-gradient tensor, describing local
stretching or compression along each axis; $\varepsilon_{xy}=\varepsilon_{yx}$
are its shear components. The tensor itself is defined at
\autoref{eq:strain-tensor}. The second equality follows from the tensor's
symmetry, which makes the two off-diagonal entries equal.

\textbf{What it does.} It reduces the four-component tensor at each pixel to
one non-negative scalar, the Frobenius norm, describing deformation
\emph{intensity} while discarding its sign and direction. Because the tensor is
symmetric, the norm counts shear twice --- once for each off-diagonal entry ---
which is what produces the factor of two.

\textbf{Source.} The Liong strain series~\cite{liong2014a,liong2014b,liong2016}.
Shreve et al.~\cite{shreve2011} establish the tensor; the Liong papers use this
four-component magnitude as the descriptor.

\textbf{As implemented.} \textbf{No.} The implementation computes the
three-term form of \autoref{eq:strain-implemented} instead. This is the
thesis's most consequential divergence and is treated in the next entry.

% -----------------------------------------------------------------------------
\subsection{Three-term strain magnitude (the implemented form)}
\label{eq:strain-implemented}
% -----------------------------------------------------------------------------

\textbf{Appears at} \autoref{sec:strain-divergence}, and again as
$\varepsilon_{\mathrm{mag}}$ at \autoref{sec:optical-strain} in Chapter~2.

\[
    \varepsilon_{\mathrm{os}}
    =\sqrt{\varepsilon_{xx}^{2}+\varepsilon_{yy}^{2}
        +\varepsilon_{xy}^{2}},
    \qquad
    \varepsilon_{xy}=\tfrac{1}{2}
        \left(\frac{\partial u}{\partial y}
             +\frac{\partial v}{\partial x}\right).
\]

Chapter~2 writes the same scalar as

\[
    \varepsilon_{\mathrm{mag}}
       =\sqrt{\varepsilon_{xx}^2+\varepsilon_{yy}^2
                    +\varepsilon_{xy}^2}.
\]

\textbf{Symbols.} As \autoref{eq:strain-frobenius}, with $u$ and $v$ the
horizontal and vertical components of the estimated flow field. In the
discrete pipeline $u$ and $v$ are displacements in pixels between selected
frames, not physical velocities.

\textbf{What it does.} The same reduction to a deformation-intensity scalar,
but counting the shear contribution once rather than twice. The consequence is
a change in relative weighting, not a removable constant: the four-term form
weights pure shear by $\sqrt{2}$ relative to this one. Independent min--max
normalisation removes a uniform scale factor but does not generally remove
this difference, because normal and shear contributions vary across pixels and
across time.

\textbf{Source.} \textbf{None --- this is a declared divergence.} STSTNet
\cite{liong2019b} cannot establish a precedent for the three-term form,
because its printed magnitude equation is typographically defective in two
independent ways. It repeats $\partial u$ in the mixed derivative where the
symmetric tensor calls for $\partial v / \partial x$ alongside $\partial u /
\partial y$, and it places the factor $\tfrac{1}{2}$ \emph{outside} the
squared derivative sum. Correcting the repeated derivative while retaining
that placement gives

\[
    \tfrac{1}{2}
       \left(\frac{\partial u}{\partial y}
            +\frac{\partial v}{\partial x}\right)^2
       =2\varepsilon_{xy}^{2},
\]

which agrees with the \emph{four}-term convention, not the three-term one.
Moving the factor inside the square would instead yield the implemented form,
but that is a further change to the printed expression. The intended
implementation cannot be settled from the published equation alone, so the
thesis declares a departure from the unambiguous four-term convention rather
than claiming to reproduce STSTNet's magnitude.

\textbf{As implemented.} This is what the code computes. Spatial derivatives
come from NumPy's \texttt{gradient} at one-pixel spacing, using central
differences in the array interior; the spacing follows Liong et
al.~\cite{liong2014a} even though the scalar reduction does not. The two
conventions were never compared experimentally here, so their practical
equivalence is not assumed. This is row 4 of the divergence ledger
(\autoref{tab:divergence-ledger}). \textbf{Neither expression should be
``corrected'' --- the difference between them is the finding.}

% -----------------------------------------------------------------------------
\subsection{SimAM: energy and gate}
\label{eq:simam}
% -----------------------------------------------------------------------------

\textbf{Appears at} \autoref{sec:simam-energy}, and at \autoref{sec:simam} in
Chapter~2. Chapter~3 writes it in terms of the intermediate $d_i$ and $s^2$:

\[
    d_i=(x_i-\mu)^2,\qquad
    s^2=\frac{1}{M-1}\sum_{i=1}^{M}d_i,\qquad
    E_i^{-1}=\frac{d_i}{4(s^2+\lambda)}+\frac{1}{2},
\]
\[
    \widetilde{x}_i=x_i\,\operatorname{sigmoid}(E_i^{-1}),
    \qquad \mu=\frac{1}{M}\sum_{i=1}^{M}x_i.
\]

Chapter~2 states the same gate with the variance written as $v$:

\[
    \mu=\frac1M\sum_i x_i,\qquad
    v=\frac1{M-1}\sum_i(x_i-\mu)^2,
\]
\[
    E_i^{-1}=\frac{(x_i-\mu)^2}{4(v+\lambda)}+\frac12,
    \qquad \widetilde x_i=x_i\operatorname{sigmoid}(E_i^{-1}).
\]

\textbf{Symbols.} $x_i$ is one activation; $\mu$ and $s^2$ (equivalently $v$)
are the mean and variance over the statistical context shared by that
channel; $M$ is the number of positions in that context; $\lambda$ is a
positive regularisation constant, fixed at $10^{-4}$ here; $E_i^{-1}$ is the
reciprocal energy, the importance score; $\widetilde{x}_i$ is the refined
activation.

\textbf{What it does.} It implements a distinctiveness prior: an activation
that differs from its channel context is scored as more important, and the
score multiplicatively rescales the feature map rather than being added to it.
The score is the closed-form minimiser of a regularised separation energy, so
the $+\tfrac{1}{2}$ follows from the reciprocal-energy expression rather than
being an arbitrary offset. The sigmoid bounds the multiplicative weight while
preserving the ordering of the scores; it does not produce a hard mask of
inactive regions. No weight or bias is learned anywhere in the module, so it
contributes \textbf{zero learnable parameters} regardless of channel count.

Note what the prior is and is not. It is a statement about activation
statistics, not a detector of facial action units or of minority classes. A
distinctive response may encode useful deformation, or noise, or an artefact.

\textbf{Source.} Yang et al.~\cite{yang2021}.

\textbf{As implemented.} Partly --- with two documented distinctions, both
deliberate.

\begin{enumerate}
	\item \textbf{The variance denominator.} The shared variance printed with
	Yang et al.'s equation~(5) divides by $M$; the reference code in their
	Figure~3 divides by $M-1$. The forms above follow the \emph{code}, as does
	this thesis, which applies Bessel's correction over $n = D \cdot H \cdot W
	- 1$ (\autoref{sec:simam-placement}).
	\item \textbf{The statistical context.} Here $M = THW$ --- whole-clip
	statistics within each channel --- against the original image
	formulation's $M = HW$ per frame. The original module produces a separate
	weight at each channel--height--width position, so its ``three-dimensional
	weights'' describe $C \times H \times W$ and not a temporal volume.
	Extending the context across frames is an additional design decision, not
	part of the source formulation, and is row 10 of the divergence ledger.
	Whole-clip statistics do not guarantee that apex frames receive the
	highest weights.
\end{enumerate}

The value $\lambda = 10^{-4}$ is likewise inherited rather than derived: Yang
et al. sweep $\lambda$ from $10^{-1}$ to $10^{-6}$ and select $10^{-4}$ for
CIFAR, repeating the selection separately for ImageNet and choosing $0.1$. It
is a dataset-selected operating point, not a universal setting justified by
the formula.

% -----------------------------------------------------------------------------
\subsection{Focal loss}
\label{eq:focal}
% -----------------------------------------------------------------------------

\textbf{Appears at} \autoref{sec:imbalance-focal}, and at
\autoref{sec:focal-loss} in Chapter~2.

\[
	\text{FL}(p_t) = -\alpha_t (1 - p_t)^{\gamma}\log p_t
\]

\textbf{Symbols.} $p_t$ is the model's predicted probability for the true
class; $\gamma$ is the focusing exponent, described by Zhao et al. as ``the
balance factor for loss''; $\alpha_t$ is a per-class multiplier, ``the weight
balance factor for samples''. With $\alpha_t = 1$ and $\gamma = 0$ the
expression reduces to ordinary cross-entropy, $-\log p_t$.

\textbf{What it does.} Two independent corrections, which are easy to
conflate and which this thesis keeps separate deliberately. The modulating
factor $(1-p_t)^{\gamma}$ re-weights by \textbf{difficulty}: for $\gamma > 0$
it shrinks the contribution of confident, correct examples relative to hard
ones. It does not inspect class frequency at all, so a difficult example from
any class retains a large contribution --- though under skew the examples it
up-weights are disproportionately from the minority classes. The separate
$\alpha_t$ multiplier re-weights by \textbf{class frequency}. Zhao et al.
treat $\gamma$ as a constant and $\alpha$ as dataset-specific, to be set by
cross-validation.

\textbf{Source.} Zhao et al.~\cite{zhaoS2021}, who bring focal loss into
micro-expression recognition ``to alleviate the inefficient model training
caused by the class imbalance in the MEs Datasets''. They adopt the
formulation developed for one-stage object detection, where the imbalance
being corrected is between foreground and background regions rather than
between emotion classes. Per the bibliography's stated policy, the original
detection paper carries no entry here and the formulation is attributed to the
reviewed paper that reports it.

\textbf{As implemented.} Partly. The difficulty term is used as written, with
$\gamma = 2$ fixed throughout, matching Zhao et al.'s reported exponent --- an
inherited setting, not a measured optimum for this corpus. The implementation
computes $p_t$ from the hard target label using log-softmax, then multiplies a
\emph{label-smoothed} cross-entropy term (\autoref{eq:label-smoothing}) by
$(1-p_t)^{\gamma}$, with the class weight applied afterwards and optionally.
The frequency term $\alpha_t$ is where this thesis departs: frequency
correction is assigned to the sampler instead, and when balanced sampling is
active a runtime guard disables the loss-side class weights, leaving focal
difficulty weighting in place. Conflating the two factors is precisely the
error recorded at \autoref{sec:imbalance-double-correction}.

% -----------------------------------------------------------------------------
\subsection{Inline quantities}
\label{eq:ch3-inline}
% -----------------------------------------------------------------------------

Three quantities in Chapter~3 are not set as display mathematics but carry the
same evidential weight.

\textbf{Flow magnitude, $\sqrt{u^2+v^2}$} (\autoref{sec:brightness-constancy}).
The scalar magnitude of the flow field, described by Zhao et
al.~\cite{zhaoS2021} alongside the component fields $u(x,y)$ and $v(x,y)$.
This thesis retains the two components as separate input channels rather than
reducing them to this magnitude, so the direction of motion is preserved.

\textbf{Effective frame rate, $6600/N$ fps} (\autoref{sec:tim-implications}).
Sampling 33 frames from a clip of $N$ frames recorded at 200~fps yields this
effective rate: \textbf{100~fps} at the median clip length $N=66$, which is
SMIC's recording rate and half what CASME~II was built to provide, and roughly
\textbf{52~fps} for the longest clip at $N=126$, below the original CASME
database's 60~fps~\cite{yan2014}. Eighty-one of the 156 clips are downsampled
by a factor of two or more. No reviewed paper reports this cost.

\textbf{True inter-frame interval, $N/(200 \times 32)$ seconds}
(\autoref{sec:tim-implications}). Magnification runs on the 33 sampled frames
at the nominal 200~fps, but those frames span the clip's full duration, so the
real interval between them is this rather than $1/200$. A band specified as
5--25~Hz therefore corresponds to roughly 2.5--12.5~Hz physically at the
median, and the discrepancy grows with clip length. The realised pass-band is
consequently clip-dependent --- neither the 5--25~Hz derived from the duration
rule of Bai et al.~\cite{bai2021} nor constant across the corpus. This is a
declared divergence; magnifying before subsampling is declared as further
work.


% =============================================================================
\section{Chapter 4 --- Methodology}
\label{sec:eq-ch4}
% =============================================================================

% -----------------------------------------------------------------------------
\subsection{Sinusoidal positional encoding}
\label{eq:posenc}
% -----------------------------------------------------------------------------

\textbf{Appears at} \autoref{sec:transformer-config}.

\[
	\mathrm{PE}(p, 2i) = \sin\!\left(\frac{p}{10000^{2i/96}}\right), \qquad \mathrm{PE}(p, 2i+1) = \cos\!\left(\frac{p}{10000^{2i/96}}\right).
\]

\textbf{Symbols.} $p$ is the position in the sequence, one of the 32 time
steps; $i$ indexes feature pairs within the embedding, so $2i$ and $2i+1$ are
the even and odd feature positions; 96 is $d_{\text{model}}$, the embedding
width, appearing here as a literal rather than a symbol because the
configuration fixes it.

\textbf{What it does.} It supplies the order information that self-attention
cannot recover on its own. Without it, attention is permutation-equivariant:
permuting the input tokens permutes the outputs correspondingly, so a shuffled
clip is indistinguishable from an ordered one (\autoref{sec:self-attention}).
Each feature pair oscillates at a different frequency, geometrically spaced by
the $10000^{2i/96}$ denominator, so the vector at each position is distinct
and positions at a fixed offset stand in a consistent linear relation.

\textbf{Source.} No corpus source. The encoding is Vaswani et al.'s, from the
original transformer, which the bibliography deliberately omits under the
policy described at the head of this appendix; Chapter~4 states it as the
standard form without attribution. The thesis reaches the architecture through
Zhang et al.'s SLSTT~\cite{zhang2022} and its Vision Transformer
foundation~\cite{dosovitskiy2021}, as traced at
\autoref{sec:transformer-review}. Note that ViT itself uses \emph{learned}
positional embeddings, not these fixed sinusoids.

\textbf{As implemented.} Yes, as written. It is a fixed buffer, not a learned
parameter, and contributes nothing to any parameter count. Two details of the
surrounding configuration are easily missed: the positional-encoding stage
applies its own \texttt{Dropout(p=0.1)} to the sequence \emph{after} the
encoding is added, which is a separate application from the encoder layers'
dropout at the same rate; and although the sinusoids could be evaluated at
arbitrary positions, the implementation uses a finite precomputed buffer, so
generalisation to unseen sequence lengths is neither exercised nor guaranteed.

% -----------------------------------------------------------------------------
\subsection{The mean-of-folds macro F1 ceiling}
\label{eq:fold-ceiling}
% -----------------------------------------------------------------------------

\textbf{Appears at} \autoref{sec:mean-of-folds-rejected}.

\[
	\frac{10 \cdot \tfrac13 + 8 \cdot \tfrac23 + 7 \cdot 1}{25} = \frac{3.333\ldots + 5.333\ldots + 7}{25} \approx 0.627.
\]

\textbf{Symbols.} The three numerator terms are the 25 leave-one-subject-out
folds grouped by how many of the three classes appear in each fold's held-out
ground truth: \textbf{10} folds contain one class, \textbf{8} contain two,
\textbf{7} contain all three. The fractions $\tfrac13$, $\tfrac23$ and $1$ are
the maximum macro F1 each group can attain. The denominator is the fold count.

\textbf{What it does.} It computes an upper bound, fixed by protocol
structure before any model is trained. A macro F1 computed \emph{within} a
fold is the unweighted mean of that fold's three per-class F1 scores, and a
class absent from that fold's ground truth contributes zero no matter what the
model predicts. A single-class fold is therefore capped at $1/3$ and a
two-class fold at $2/3$. Averaging across folds averages seventeen sub-unity
ceilings into every score, bounding the mean-of-folds macro F1 at
approximately \textbf{0.627} --- a property of which subjects were held out,
identical for all twelve configurations.

Pooling avoids the ceiling entirely: accumulating $TP$, $FP$ and $FN$ across
all 25 folds and all 156 clips before computing any per-class F1 means a class
locally absent from one fold no longer forces a zero into its contribution.
That is the construction of \autoref{eq:uf1}.

\textbf{Source.} This thesis's own arithmetic over the fold composition of
\autoref{fig:fold-composition}. It is not a published result.

\textbf{As implemented.} Not applicable --- this is an argument for rejecting
a metric, not a component. Its consequence is directly implemented: pooled
macro F1 is the primary reported metric, and the mean-of-folds quantity stored
under \texttt{macro\_f1} in the run artefacts is not. A metric with a hard
ceiling below 0.7 cannot serve as a headline result. The same fold composition
biases mean-of-folds \emph{accuracy} in the same direction without imposing a
hard ceiling; for the full configuration it exceeds pooled accuracy by
\textbf{6.3 points}.

% -----------------------------------------------------------------------------
\subsection{Inline quantities}
\label{eq:ch4-inline}
% -----------------------------------------------------------------------------

\textbf{Sampler weight, $1/n_c$} (\autoref{sec:sampler-config}). Every
training example is assigned this weight by the
\texttt{WeightedRandomSampler}, where $n_c$ is the number of training examples
of that example's class \emph{within the current fold's training split}.
Sampling proceeds with \texttt{replacement=True}, drawing as many indices per
epoch as the split contains, which makes each represented class equally likely
in expectation while leaving individual batches uneven. This changes class
exposure during training, not the number of distinct recorded expressions.
It is also what makes the class-weighting rule well founded rather than merely
convenient: because the skew correction genuinely happens here, disabling its
loss-side equivalent avoids double-counting rather than removing the
correction (\autoref{eq:focal}).

\textbf{Parameter counts} (\autoref{sec:parameter-counts}). Four components
combine to give any configuration's total: the convolutional backbone
(\texttt{ThreeStreamCNNBackbone}, all three streams) \textbf{14,544}; the
transformer encoder (\texttt{SLSTTTransformer}) \textbf{348,736};
\texttt{RawPatchEmbedding}, the no-CNN fallback, \textbf{4,704}; and the
classifier head \textbf{483}. Their sums give the four distinct totals of
\autoref{tab:parameter-totals}. Neither SimAM nor EVM appears as a row: SimAM
contributes zero parameters by construction (\autoref{eq:simam}), and
magnification is a preprocessing step rather than a network component. Of the
transformer's total, 192 belong to the final \texttt{LayerNorm(96)} applied
after its four encoder layers.


% =============================================================================
\section{Chapter 2 --- Background}
\label{sec:eq-ch2}
% =============================================================================

% -----------------------------------------------------------------------------
\subsection{Eulerian magnification: the small-motion approximation}
\label{eq:evm-taylor}
% -----------------------------------------------------------------------------

\textbf{Appears at} \autoref{sec:eulerian-idea}.

\[
    I(x,t)=f(x+\delta(t))
       \approx f(x)+\delta(t)\frac{\partial f}{\partial x}.
\]

\textbf{Symbols.} $f$ is the image intensity profile; $x$ is spatial
position; $\delta(t)$ is a small time-varying displacement; $I(x,t)$ is the
observed intensity at a \emph{fixed} location $x$ at time $t$.

\textbf{What it does.} It is the first-order Taylor expansion that licenses
the entire Eulerian approach. Eulerian magnification examines intensity
changes at fixed image locations rather than explicitly tracking moving
points, and this approximation is what connects the two: for sufficiently
small displacement, the intensity change observed at a fixed pixel is
proportional to the displacement times the local spatial gradient. Motion can
therefore be manipulated by manipulating intensity.

\textbf{Source.} Bai et al.~\cite{bai2021}, describing the amplitude-based
mechanism. The technique originates with Wu et al., who carry no bibliography
entry under the policy described above and are attributed through the reviewed
paper that reports them.

\textbf{As implemented.} Yes --- but note that the approximation holds only
for small $\delta(t)$. It does not imply that arbitrary intensity changes
represent facial movement, nor that larger amplification always improves
recognition.

% -----------------------------------------------------------------------------
\subsection{Eulerian magnification: the amplified signal}
\label{eq:evm-amplified}
% -----------------------------------------------------------------------------

\textbf{Appears at} \autoref{sec:eulerian-idea}.

\[
    \widetilde I(x,t)\approx
       f(x)+(1+\alpha)\delta(t)\frac{\partial f}{\partial x}.
\]

\textbf{Symbols.} As above, with $\alpha$ the magnification factor and
$\widetilde I$ the reconstructed signal.

\textbf{What it does.} Adding $\alpha$ times the band-passed motion-related
variation back to the original produces a signal that, under the same
first-order approximation, resembles what would have been observed had the
displacement itself been $(1+\alpha)$ times larger. The mechanism is entirely
photometric; nothing is tracked or warped.

\textbf{Source.} Bai et al.~\cite{bai2021}, as above.

\textbf{As implemented.} \textbf{Not in the temporal filter.} The magnitude
relation above is implemented, but the band-pass stage that isolates
$\delta(t)$ differs: Bai et al. describe a Butterworth filter, whereas this
implementation masks Fourier coefficients outside the selected band and
applies the inverse transform (\autoref{sec:temporal-filtering}). The spatial
decomposition uses a Laplacian pyramid with the separable blur kernel
$[1,4,6,4,1]/16$ and factor-two downsampling. A second and larger divergence
concerns \emph{when} magnification is applied: it runs after subsampling, which
makes the realised pass-band clip-dependent (\autoref{eq:ch3-inline}).

% -----------------------------------------------------------------------------
\subsection{Brightness constancy}
\label{eq:brightness}
% -----------------------------------------------------------------------------

\textbf{Appears at} \autoref{sec:optical-flow}.

\[
    I(x,y,t)=I(x+u\Delta t,y+v\Delta t,t+\Delta t).
\]

\textbf{Symbols.} $I$ is image intensity; $(x,y)$ is spatial position; $t$ is
time; $u$ and $v$ are the horizontal and vertical components of local
velocity, for the continuous derivation.

\textbf{What it does.} It states the assumption on which optical flow rests:
a moving point retains its intensity as it moves. Everything estimated as
``motion'' downstream is motion only insofar as this holds.

\textbf{Source.} OFF-ApexNet, Liong et al.~\cite{liong2019a}.

\textbf{As implemented.} Yes, as the assumption underlying the Farneb\"ack
estimator used here --- with two caveats the thesis states rather than
assumes away. In the discrete pipeline $u$ and $v$ denote displacement in
pixels between selected frames rather than physical velocity per second. And
illumination changes or large inter-frame displacement can violate the
approximation outright, which the clip-dependent effective frame rate
(\autoref{eq:ch3-inline}) makes a live concern for the longer clips.

% -----------------------------------------------------------------------------
\subsection{The optical flow constraint}
\label{eq:flow-constraint}
% -----------------------------------------------------------------------------

\textbf{Appears at} \autoref{sec:optical-flow}.

\[
    I_xu+I_yv+I_t=0.
\]

\textbf{Symbols.} $I_x$, $I_y$ and $I_t$ are the partial derivatives of
intensity with respect to $x$, $y$ and $t$; $u$ and $v$ are as above.

\textbf{What it does.} It is the first-order expansion of brightness
constancy, and it is \textbf{one equation in two unknowns}. This
underdetermination is the \textbf{aperture problem}: a local edge constrains
motion normal to the edge, along the intensity gradient, but leaves motion
parallel to the edge entirely undetermined. Every practical estimator must
therefore add something --- a neighbourhood assumption, a local motion model,
or explicit regularisation. Flow fields are consequently never pure
measurements; they carry the assumptions of whichever estimator produced them.

\textbf{Source.} Liong et al.~\cite{liong2019a}.

\textbf{As implemented.} Yes. The constraint is resolved by Farneb\"ack's
polynomial expansion, following the estimator choice reported by Zhao et
al.~\cite{zhaoS2021}, with OpenCV settings: pyramid scale 0.5, three levels,
window size 15, three iterations, polynomial neighbourhood 5, polynomial
smoothing 1.2, flags 0. Flow is computed for the 32 adjacent pairs among the
33 sampled frames.

% -----------------------------------------------------------------------------
\subsection{The symmetric displacement-gradient tensor}
\label{eq:strain-tensor}
% -----------------------------------------------------------------------------

\textbf{Appears at} \autoref{sec:optical-strain}. Its normal and shear
components are also named at \autoref{sec:strain-tensor} in Chapter~3.

\[
    \varepsilon=\tfrac12\left(\nabla\mathbf{u}
                         +(\nabla\mathbf{u})^{\mathsf T}\right),
    \qquad
    \varepsilon_{xx}=\frac{\partial u}{\partial x},\quad
    \varepsilon_{yy}=\frac{\partial v}{\partial y},
\]
\[
    \varepsilon_{xy}=\varepsilon_{yx}
       =\tfrac12\left(\frac{\partial u}{\partial y}
                    +\frac{\partial v}{\partial x}\right).
\]

\textbf{Symbols.} $\mathbf{u}=[u,v]^{\mathsf T}$ is the displacement field;
$\nabla\mathbf{u}$ is its gradient; the superscript $\mathsf{T}$ denotes
transpose. Symmetrising the gradient makes $\varepsilon_{xy}=\varepsilon_{yx}$
by construction.

\textbf{What it does.} It converts displacement into \emph{deformation}.
Flow describes how much things moved; strain describes how neighbouring
displacements \emph{differ}. The normal components describe local stretching or
compression along each axis, and the off-diagonal terms describe shear. The
consequence that motivates using it at all: a spatially uniform translation
has zero displacement gradient and therefore zero strain, so rigid head
movement is suppressed while differential skin motion survives.

\textbf{Source.} The small-strain approximation used for optical strain,
Shreve et al.~\cite{shreve2011}.

\textbf{As implemented.} Yes, with the discretisation as a stated modelling
choice rather than a neutral detail, since it sets the scale at which local
variation is measured. Shreve et al. approximate the spatial derivatives by
central differences with offsets of approximately two to three pixels; Liong
et al.~\cite{liong2014a} use one-pixel offsets, and this implementation
follows the latter via NumPy's \texttt{gradient}. Two limits are declared
rather than argued away: zero strain under uniform translation does not imply
immunity to head \emph{rotation}, projection effects or flow-estimation error,
and differentiation may amplify noise. The reduction of this tensor to a
scalar is where the thesis diverges from the literature --- see
\autoref{eq:strain-implemented}.

% -----------------------------------------------------------------------------
\subsection{Convolution parameter count}
\label{eq:conv-params}
% -----------------------------------------------------------------------------

\textbf{Appears at} \autoref{sec:convolution}.

\[
    N_{\mathrm{params}}=C_{\mathrm{out}}
       (C_{\mathrm{in}}k_Hk_W+1).
\]

\textbf{Symbols.} $C_{\mathrm{in}}$ and $C_{\mathrm{out}}$ are the input and
output channel counts; $k_H$ and $k_W$ are the kernel height and width; the
$+1$ is one bias per output channel.

\textbf{What it does.} It makes weight sharing quantitative. Each output
channel's kernel spans all input channels over its local window, and the same
kernel is applied at every position --- so the parameter count is
\textbf{independent of image dimensions}. That independence is the point of
convolution: no separate detector is learned for each image location. What
does scale with image size is computation and activation memory, since both
grow with the number of positions processed.

\textbf{Source.} Standard; no corpus source and no citation.

\textbf{As implemented.} Yes. The implemented streams use $(1,3,3)$ kernels
with $(0,1,1)$ padding followed by $(1,2,2)$ max-pooling, so the convolutions
preserve temporal length and learn purely spatial filters while pooling halves
height and width. Because $k_T = 1$, \textbf{no temporal convolution is
performed} --- the use of \texttt{Conv3d} does not by itself establish
otherwise. The complete stream can still depend on multiple frames through
BatchNorm during training and through SimAM whenever it is enabled, both of
which use statistics spanning time.

% -----------------------------------------------------------------------------
\subsection{Batch normalisation}
\label{eq:batchnorm}
% -----------------------------------------------------------------------------

\textbf{Appears at} \autoref{sec:normalisation}.

\[
    \widehat x=\frac{x-\mu_c}{\sqrt{\sigma_c^2+\epsilon}},
    \qquad y=\gamma_c\widehat x+\beta_c.
\]

\textbf{Symbols.} $\mu_c$ and $\sigma_c^2$ are the mean and variance for
channel $c$; $\epsilon$ is a small constant preventing division by zero;
$\gamma_c$ and $\beta_c$ are the learned per-channel scale and shift;
$\widehat x$ is the standardised activation and $y$ the output.

\textbf{What it does.} It rescales each channel's activations to a
standardised distribution and then lets the network learn to undo or reshape
that standardisation through $\gamma_c$ and $\beta_c$ --- so normalisation is
imposed but not forced to be retained.

\textbf{Source.} \textbf{None.} Batch normalisation is Ioffe and Szegedy's,
but no paper in this thesis's corpus introduces or examines it, so it is
presented as standard practice and carries no citation. Adding one would
require a new bibliography entry, which the citation policy of this thesis
excludes.

\textbf{As implemented.} Yes. \texttt{BatchNorm3d} computes training
statistics over $(B,T,H,W)$ for each channel, and evaluation uses running
estimates accumulated during training. These activation statistics are a
distinct thing from the per-sample input standardisation performed by the
dataset loader, and the two should not be conflated. Related regularisers in
the same configuration: \texttt{Dropout3d} removes entire feature channels for
a sample rather than masking individual spatial values, and layer
normalisation --- which standardises a token's feature dimension independently
of other samples --- appears inside the transformer and in the classifier
head.

% -----------------------------------------------------------------------------
\subsection{Scaled dot-product attention}
\label{eq:attention}
% -----------------------------------------------------------------------------

\textbf{Appears at} \autoref{sec:self-attention}.

\[
    \operatorname{Attention}(Q,K,V)
      =\operatorname{softmax}\left(\frac{QK^{\mathsf T}}{\sqrt{d_k}}\right)V.
\]

\textbf{Symbols.} $Q$, $K$ and $V$ are the query, key and value matrices,
each a projection of the input tokens; $d_k$ is the key dimension; the
softmax is taken row-wise.

\textbf{What it does.} Queries and keys determine \emph{which} values
contribute to each output, and by how much: the dot product $QK^{\mathsf T}$
scores every token pair, the softmax turns those scores into weights summing
to one, and the result mixes the values accordingly. The $\sqrt{d_k}$ divisor
moderates dot-product magnitude as the key dimension grows, keeping the
softmax out of its saturated regime. Multiple heads learn projections into
different subspaces, whose outputs are concatenated and linearly mixed.

The property that matters most here is negative: without position
information, attention is \textbf{permutation-equivariant}. Permuting the
input tokens permutes the outputs correspondingly, so the operation cannot
identify temporal order on its own. That is what
\autoref{eq:posenc} exists to supply.

\textbf{Source.} Vaswani et al.'s transformer, attributed through Dosovitskiy
et al.~\cite{dosovitskiy2021}, whose Vision Transformer is the foundation this
thesis reaches the architecture through, together with Zhang et al.'s
SLSTT~\cite{zhang2022} (\autoref{sec:transformer-review}). The original
transformer paper carries no bibliography entry under the stated policy.

\textbf{As implemented.} Yes, within \texttt{nn.TransformerEncoder}, but over
a different axis than in the published SLSTT. Here attention operates over
\textbf{32 time steps} of 96-dimensional features, not over spatial patches:
ViT splits an image into patches and relates them spatially, whereas the
sequence here is produced by collapsing both spatial dimensions with
\texttt{AdaptiveAvgPool3d((32,1,1))}. The configuration is $d_{\text{model}} =
96$, $n_{\text{head}} = 8$, four encoder layers, feed-forward dimension 256,
dropout 0.1, pre-norm ordering, GELU activation.

% -----------------------------------------------------------------------------
\subsection{Label smoothing}
\label{eq:label-smoothing}
% -----------------------------------------------------------------------------

\textbf{Appears at} \autoref{sec:label-smoothing}.

\[
    L_{\mathrm{smooth}}
      =-(1-\epsilon)\log p_t
        -\frac{\epsilon}{C}\sum_{c=1}^{C}\log p_c.
\]

\textbf{Symbols.} $\epsilon$ is the smoothing strength; $C$ is the number of
classes; $p_t$ is the predicted probability of the true class and $p_c$ of
class $c$. The target it corresponds to is
$q_c=(1-\epsilon)\mathbf{1}[c=t]+\epsilon/C$, the one-hot vector blended with
uniform mass.

\textbf{What it does.} It removes the incentive to drive the true-class
probability to exactly one. Because a share $\epsilon$ of the target mass sits
on every class, the loss is minimised by a confident but not saturated
prediction. The effect is on \emph{calibration}, not on class balance: it
discourages extreme confidence uniformly and does not preferentially weight a
rare class.

\textbf{Source.} Named, but not examined, by Dosovitskiy et
al.~\cite{dosovitskiy2021}, who list it among the regularisation parameters
they tune. \textbf{No reviewed micro-expression paper uses it, discusses it,
or reports what it does under class skew.}

\textbf{As implemented.} Yes, with $\epsilon = 0.05$, applied \emph{before}
focal modulation --- the focal factor of \autoref{eq:focal} still uses the
unsmoothed true-class probability. Neither the smoothing strength nor the
focal exponent is varied across the architectural comparisons. Because the
corpus supports the choice neither way, it is declared as a choice rather than
defended as a finding, and it appears as such in the divergence ledger.

% -----------------------------------------------------------------------------
\subsection{Precision, recall and per-class F1}
\label{eq:prf}
% -----------------------------------------------------------------------------

\textbf{Appears at} \autoref{sec:confusion-matrix}. The $F1_c$ term also
appears within \autoref{eq:uf1}.

\[
    P_c=\frac{TP_c}{TP_c+FP_c},\qquad
    R_c=\frac{TP_c}{TP_c+FN_c},\qquad
    F1_c=\frac{2TP_c}{2TP_c+FP_c+FN_c}.
\]

\textbf{Symbols.} In a confusion matrix, entry $(i,j)$ counts examples whose
true class is $i$ and predicted class is $j$; the diagonal holds correct
predictions. For class $c$: $TP_c$ counts correct predictions of $c$, $FP_c$
counts incorrect predictions of $c$, and $FN_c$ counts missed examples of $c$.

\textbf{What it does.} Precision penalises false alarms and recall penalises
misses --- each is trivially maximised at the other's expense, which is why
neither is reported alone. F1 is their harmonic mean, so a high value requires
both to be reasonably high; the form above is the algebraic simplification
that avoids computing $P_c$ and $R_c$ separately.

\textbf{Source.} Standard; no corpus source and no citation.

\textbf{As implemented.} Yes, with one convention made explicit because it is
load-bearing under this fold structure: a degenerate denominator, which occurs
whenever a class is absent from a fold, is assigned \textbf{zero} rather than
the class being omitted from the score. That convention is exactly what
produces the ceiling of \autoref{eq:fold-ceiling}, and exactly what pooling
across folds avoids.

% -----------------------------------------------------------------------------
\subsection{Macro-averaged F1}
\label{eq:macro-f1}
% -----------------------------------------------------------------------------

\textbf{Appears at} \autoref{sec:macro-micro}.

\[
    F1_{\mathrm{macro}}=\frac1C\sum_{c=1}^{C}F1_c.
\]

\textbf{Symbols.} $C$ is the number of classes and $F1_c$ is the per-class F1
of \autoref{eq:prf}.

\textbf{What it does.} It gives every class equal weight regardless of how
many examples it has, which is the whole reason for preferring it here.
Contrast micro F1, which first sums $TP$, $FP$ and $FN$ \emph{across classes}:
in single-label multiclass classification each error supplies one false
positive and one false negative, so micro F1 equals accuracy exactly, and can
therefore be dominated by a frequent class. Macro F1 and accuracy describe
different aspects of performance, and reporting both does not remove the need
to inspect per-class errors directly.

\textbf{Source.} Standard. See et al.~\cite{see2019} define the \emph{pooled}
variant, MEGC's UF1, which is the construction this thesis uses
(\autoref{eq:uf1}).

\textbf{As implemented.} Yes, and it is the primary metric --- but only in
its pooled form. Pooling the confusion matrices across folds before computing
any per-class F1 is a genuinely different estimator from computing macro F1
within each fold and averaging the results (\autoref{sec:pooled-vs-fold}). The
distinction is the subject of \autoref{eq:fold-ceiling}, and no run artefact
stores the pooled quantity under any key: it is derived afterwards by the
reporting scripts under \texttt{tools/}.

%
% It must come after \appendix. Including it as a numbered chapter would shift
% the chapter numbering that chapter3_shortened.tex's divergence ledger depends
% on (it cites ten section numbers literally: 3.1.2, 3.1.3, 3.2.6, 3.4.3,
% 3.4.7, 3.5.7, 3.6.6, 3.7.7, 3.8.6, 3.9.7).
%
% CITATION POLICY: this file adds NO new bibliography entries. Works cited only
% inside the reviewed papers (Vaswani, Lin et al., Ioffe & Szegedy, Wu et al.)
% are attributed in prose to the corpus paper that reports them, per the policy
% stated in the header of bib/thesis.bib. Adding an entry would renumber the
% whole bibliography, since tools/bibliography.py numbers by first appearance.
%
% All labels in this file are prefixed eq: and collide with nothing.
% =============================================================================

\chapter{Mathematical Formulations}
\label{ch:formulas}

This appendix collects every display formula that appears in the body of the
thesis, in one place, with a consistent treatment: the expression as it is
written in its own chapter, the section it belongs to, the meaning of each
symbol, the mechanism it implements, the work it comes from, and whether this
thesis implements it as written.

\section*{How to read this appendix}
\addcontentsline{toc}{section}{How to read this appendix}

\textbf{Ordering.} Sequential by chapter, in the order the mathematics is
introduced: Chapter~3 (\autoref{sec:eq-ch3}), Chapter~4
(\autoref{sec:eq-ch4}), then Chapter~2 (\autoref{sec:eq-ch2}), which carries
the bulk of the derivations. Chapters~1, 5 and~6 contain no display
mathematics.

\textbf{Reproduction.} Every expression is reproduced exactly as it is
typeset in its own chapter --- same symbols, same subscripts, same
conventions. Where two chapters state the same relation in different notation,
both forms are given inside the one entry rather than being silently
reconciled.

\textbf{Each formula appears once.} Chapter~2 develops the mechanism and
Chapter~3 reports the literature's version and any divergence from it, so
several relations occur in both. Those are presented once, at the location
listed first in the ordering above, with a cross-reference recording every
section that states them. No entry is duplicated.

\textbf{Attribution.} Citations use only keys already present in
\texttt{bib/thesis.bib}. Several formulas originate in works that the
bibliography deliberately omits, because this thesis's corpus is the set of
reviewed micro-expression papers rather than the full ancestry of every
technique they use. Those are attributed in prose to the reviewed paper that
brings them to this task --- focal loss through Zhao et
al.~\cite{zhaoS2021}, self-attention and the transformer through Dosovitskiy
et al.~\cite{dosovitskiy2021} and Zhang et al.~\cite{zhang2022}, Eulerian
magnification through Bai et al.~\cite{bai2021}. Formulas that are standard
textbook material with no corpus source, such as batch normalisation and the
convolution parameter count, carry no citation and are marked as such.

\textbf{Implementation status.} Each entry closes by stating whether the
thesis implements the formula as written. Three do not, and those are declared
divergences recorded in the divergence ledger of \autoref{tab:divergence-ledger}
rather than defects to be corrected here.

\begin{table}[htbp]
	\centering
	\footnotesize
	\begin{tabular}{@{}l l l l@{}}
		\hline
		\textbf{Entry} & \textbf{Quantity} & \textbf{Attributed to} & \textbf{As written?} \\
		\hline
		\multicolumn{4}{@{}l}{\textit{Chapter 3 --- Literature Review}} \\
		\autoref{eq:uf1}           & Per-class F1 and UF1          & \cite{see2019}                        & yes \\
		\autoref{eq:strain-frobenius} & Four-term strain magnitude & \cite{liong2014a,liong2014b,liong2016} & \textbf{no} \\
		\autoref{eq:strain-implemented} & Three-term strain magnitude & declared divergence                & --- \\
		\autoref{eq:simam}         & SimAM energy and gate         & \cite{yang2021}                       & partly \\
		\autoref{eq:focal}         & Focal loss                    & \cite{zhaoS2021}                      & partly \\
		\multicolumn{4}{@{}l}{\textit{Chapter 4 --- Methodology}} \\
		\autoref{eq:posenc}        & Sinusoidal positional encoding & no corpus source                     & yes \\
		\autoref{eq:fold-ceiling}  & Mean-of-folds macro F1 ceiling & this thesis                          & n/a \\
		\multicolumn{4}{@{}l}{\textit{Chapter 2 --- Background}} \\
		\autoref{eq:evm-taylor}    & Small-motion approximation    & \cite{bai2021}                        & yes \\
		\autoref{eq:evm-amplified} & Amplified signal              & \cite{bai2021}                        & \textbf{no} \\
		\autoref{eq:brightness}    & Brightness constancy          & \cite{liong2019a}                     & yes \\
		\autoref{eq:flow-constraint} & Optical flow constraint     & \cite{liong2019a}                     & yes \\
		\autoref{eq:strain-tensor} & Symmetric displacement gradient & \cite{shreve2011}                   & yes \\
		\autoref{eq:conv-params}   & Convolution parameter count   & standard                              & yes \\
		\autoref{eq:batchnorm}     & Batch normalisation           & standard, no corpus source            & yes \\
		\autoref{eq:attention}     & Scaled dot-product attention  & via \cite{dosovitskiy2021}            & yes \\
		\autoref{eq:label-smoothing} & Label smoothing             & named by \cite{dosovitskiy2021}       & yes \\
		\autoref{eq:prf}           & Precision, recall, F1         & standard                              & yes \\
		\autoref{eq:macro-f1}      & Macro-averaged F1             & standard; \cite{see2019} pooled       & yes \\
		\hline
	\end{tabular}
	\caption[Index of formulations]{Index of the formulations collected in this
	appendix, with the work each is attributed to and whether this thesis
	implements it as written.}
	\label{tab:formula-index}
\end{table}


% =============================================================================
\section{Chapter 3 --- Literature Review}
\label{sec:eq-ch3}
% =============================================================================

% -----------------------------------------------------------------------------
\subsection{Per-class F1 and Unweighted F1}
\label{eq:uf1}
% -----------------------------------------------------------------------------

\textbf{Appears at} \autoref{sec:megc-metrics}.

\[
	F1_c = \frac{2 \cdot TP_c}{2 \cdot TP_c + FP_c + FN_c}, \qquad \text{UF1} = \frac{1}{C}\sum_c F1_c
\]

\textbf{Symbols.} $TP_c$, $FP_c$ and $FN_c$ are the true positives, false
positives and false negatives for class $c$; $C$ is the number of classes,
three here; $F1_c$ is the harmonic mean of precision and recall for class $c$,
written in the algebraically equivalent form that avoids computing either
separately.

\textbf{What it does.} UF1 scores a classifier without letting a frequent
class dominate, by computing a score per class and averaging the scores with
equal weight rather than averaging over examples. The construction that
matters is not visible in the algebra: the counts $TP_c$, $FP_c$ and $FN_c$
are \emph{accumulated across all $k$ folds before any $F1_c$ is computed}.
That pooling is what distinguishes UF1 from the average of per-fold macro F1
scores, which is a different estimator with a different expectation
(\autoref{sec:pooled-vs-fold}, and \autoref{eq:fold-ceiling} for what it costs).

\textbf{Source.} The metric convention of the Second Facial Micro-Expression
Grand Challenge, See et al.~\cite{see2019}, which mandates UF1 and Unweighted
Average Recall in place of plain accuracy on the explicit grounds of class
imbalance.

\textbf{As implemented.} Yes, and it is the thesis's primary metric. Pooled
macro F1 as computed here is identical in construction to MEGC's UF1. The
thesis adopts the challenge's metric definitions while declining its composite
training regime.

% -----------------------------------------------------------------------------
\subsection{Four-term strain magnitude (the published convention)}
\label{eq:strain-frobenius}
% -----------------------------------------------------------------------------

\textbf{Appears at} \autoref{sec:strain-tensor}.

\[
    \varepsilon_{\mathrm{F}}
    =\sqrt{\varepsilon_{xx}^{2}+\varepsilon_{yy}^{2}
        +\varepsilon_{xy}^{2}+\varepsilon_{yx}^{2}}
    =\sqrt{\varepsilon_{xx}^{2}+\varepsilon_{yy}^{2}
        +2\varepsilon_{xy}^{2}}.
\]

\textbf{Symbols.} $\varepsilon_{xx}$ and $\varepsilon_{yy}$ are the normal
components of the symmetric displacement-gradient tensor, describing local
stretching or compression along each axis; $\varepsilon_{xy}=\varepsilon_{yx}$
are its shear components. The tensor itself is defined at
\autoref{eq:strain-tensor}. The second equality follows from the tensor's
symmetry, which makes the two off-diagonal entries equal.

\textbf{What it does.} It reduces the four-component tensor at each pixel to
one non-negative scalar, the Frobenius norm, describing deformation
\emph{intensity} while discarding its sign and direction. Because the tensor is
symmetric, the norm counts shear twice --- once for each off-diagonal entry ---
which is what produces the factor of two.

\textbf{Source.} The Liong strain series~\cite{liong2014a,liong2014b,liong2016}.
Shreve et al.~\cite{shreve2011} establish the tensor; the Liong papers use this
four-component magnitude as the descriptor.

\textbf{As implemented.} \textbf{No.} The implementation computes the
three-term form of \autoref{eq:strain-implemented} instead. This is the
thesis's most consequential divergence and is treated in the next entry.

% -----------------------------------------------------------------------------
\subsection{Three-term strain magnitude (the implemented form)}
\label{eq:strain-implemented}
% -----------------------------------------------------------------------------

\textbf{Appears at} \autoref{sec:strain-divergence}, and again as
$\varepsilon_{\mathrm{mag}}$ at \autoref{sec:optical-strain} in Chapter~2.

\[
    \varepsilon_{\mathrm{os}}
    =\sqrt{\varepsilon_{xx}^{2}+\varepsilon_{yy}^{2}
        +\varepsilon_{xy}^{2}},
    \qquad
    \varepsilon_{xy}=\tfrac{1}{2}
        \left(\frac{\partial u}{\partial y}
             +\frac{\partial v}{\partial x}\right).
\]

Chapter~2 writes the same scalar as

\[
    \varepsilon_{\mathrm{mag}}
       =\sqrt{\varepsilon_{xx}^2+\varepsilon_{yy}^2
                    +\varepsilon_{xy}^2}.
\]

\textbf{Symbols.} As \autoref{eq:strain-frobenius}, with $u$ and $v$ the
horizontal and vertical components of the estimated flow field. In the
discrete pipeline $u$ and $v$ are displacements in pixels between selected
frames, not physical velocities.

\textbf{What it does.} The same reduction to a deformation-intensity scalar,
but counting the shear contribution once rather than twice. The consequence is
a change in relative weighting, not a removable constant: the four-term form
weights pure shear by $\sqrt{2}$ relative to this one. Independent min--max
normalisation removes a uniform scale factor but does not generally remove
this difference, because normal and shear contributions vary across pixels and
across time.

\textbf{Source.} \textbf{None --- this is a declared divergence.} STSTNet
\cite{liong2019b} cannot establish a precedent for the three-term form,
because its printed magnitude equation is typographically defective in two
independent ways. It repeats $\partial u$ in the mixed derivative where the
symmetric tensor calls for $\partial v / \partial x$ alongside $\partial u /
\partial y$, and it places the factor $\tfrac{1}{2}$ \emph{outside} the
squared derivative sum. Correcting the repeated derivative while retaining
that placement gives

\[
    \tfrac{1}{2}
       \left(\frac{\partial u}{\partial y}
            +\frac{\partial v}{\partial x}\right)^2
       =2\varepsilon_{xy}^{2},
\]

which agrees with the \emph{four}-term convention, not the three-term one.
Moving the factor inside the square would instead yield the implemented form,
but that is a further change to the printed expression. The intended
implementation cannot be settled from the published equation alone, so the
thesis declares a departure from the unambiguous four-term convention rather
than claiming to reproduce STSTNet's magnitude.

\textbf{As implemented.} This is what the code computes. Spatial derivatives
come from NumPy's \texttt{gradient} at one-pixel spacing, using central
differences in the array interior; the spacing follows Liong et
al.~\cite{liong2014a} even though the scalar reduction does not. The two
conventions were never compared experimentally here, so their practical
equivalence is not assumed. This is row 4 of the divergence ledger
(\autoref{tab:divergence-ledger}). \textbf{Neither expression should be
``corrected'' --- the difference between them is the finding.}

% -----------------------------------------------------------------------------
\subsection{SimAM: energy and gate}
\label{eq:simam}
% -----------------------------------------------------------------------------

\textbf{Appears at} \autoref{sec:simam-energy}, and at \autoref{sec:simam} in
Chapter~2. Chapter~3 writes it in terms of the intermediate $d_i$ and $s^2$:

\[
    d_i=(x_i-\mu)^2,\qquad
    s^2=\frac{1}{M-1}\sum_{i=1}^{M}d_i,\qquad
    E_i^{-1}=\frac{d_i}{4(s^2+\lambda)}+\frac{1}{2},
\]
\[
    \widetilde{x}_i=x_i\,\operatorname{sigmoid}(E_i^{-1}),
    \qquad \mu=\frac{1}{M}\sum_{i=1}^{M}x_i.
\]

Chapter~2 states the same gate with the variance written as $v$:

\[
    \mu=\frac1M\sum_i x_i,\qquad
    v=\frac1{M-1}\sum_i(x_i-\mu)^2,
\]
\[
    E_i^{-1}=\frac{(x_i-\mu)^2}{4(v+\lambda)}+\frac12,
    \qquad \widetilde x_i=x_i\operatorname{sigmoid}(E_i^{-1}).
\]

\textbf{Symbols.} $x_i$ is one activation; $\mu$ and $s^2$ (equivalently $v$)
are the mean and variance over the statistical context shared by that
channel; $M$ is the number of positions in that context; $\lambda$ is a
positive regularisation constant, fixed at $10^{-4}$ here; $E_i^{-1}$ is the
reciprocal energy, the importance score; $\widetilde{x}_i$ is the refined
activation.

\textbf{What it does.} It implements a distinctiveness prior: an activation
that differs from its channel context is scored as more important, and the
score multiplicatively rescales the feature map rather than being added to it.
The score is the closed-form minimiser of a regularised separation energy, so
the $+\tfrac{1}{2}$ follows from the reciprocal-energy expression rather than
being an arbitrary offset. The sigmoid bounds the multiplicative weight while
preserving the ordering of the scores; it does not produce a hard mask of
inactive regions. No weight or bias is learned anywhere in the module, so it
contributes \textbf{zero learnable parameters} regardless of channel count.

Note what the prior is and is not. It is a statement about activation
statistics, not a detector of facial action units or of minority classes. A
distinctive response may encode useful deformation, or noise, or an artefact.

\textbf{Source.} Yang et al.~\cite{yang2021}.

\textbf{As implemented.} Partly --- with two documented distinctions, both
deliberate.

\begin{enumerate}
	\item \textbf{The variance denominator.} The shared variance printed with
	Yang et al.'s equation~(5) divides by $M$; the reference code in their
	Figure~3 divides by $M-1$. The forms above follow the \emph{code}, as does
	this thesis, which applies Bessel's correction over $n = D \cdot H \cdot W
	- 1$ (\autoref{sec:simam-placement}).
	\item \textbf{The statistical context.} Here $M = THW$ --- whole-clip
	statistics within each channel --- against the original image
	formulation's $M = HW$ per frame. The original module produces a separate
	weight at each channel--height--width position, so its ``three-dimensional
	weights'' describe $C \times H \times W$ and not a temporal volume.
	Extending the context across frames is an additional design decision, not
	part of the source formulation, and is row 10 of the divergence ledger.
	Whole-clip statistics do not guarantee that apex frames receive the
	highest weights.
\end{enumerate}

The value $\lambda = 10^{-4}$ is likewise inherited rather than derived: Yang
et al. sweep $\lambda$ from $10^{-1}$ to $10^{-6}$ and select $10^{-4}$ for
CIFAR, repeating the selection separately for ImageNet and choosing $0.1$. It
is a dataset-selected operating point, not a universal setting justified by
the formula.

% -----------------------------------------------------------------------------
\subsection{Focal loss}
\label{eq:focal}
% -----------------------------------------------------------------------------

\textbf{Appears at} \autoref{sec:imbalance-focal}, and at
\autoref{sec:focal-loss} in Chapter~2.

\[
	\text{FL}(p_t) = -\alpha_t (1 - p_t)^{\gamma}\log p_t
\]

\textbf{Symbols.} $p_t$ is the model's predicted probability for the true
class; $\gamma$ is the focusing exponent, described by Zhao et al. as ``the
balance factor for loss''; $\alpha_t$ is a per-class multiplier, ``the weight
balance factor for samples''. With $\alpha_t = 1$ and $\gamma = 0$ the
expression reduces to ordinary cross-entropy, $-\log p_t$.

\textbf{What it does.} Two independent corrections, which are easy to
conflate and which this thesis keeps separate deliberately. The modulating
factor $(1-p_t)^{\gamma}$ re-weights by \textbf{difficulty}: for $\gamma > 0$
it shrinks the contribution of confident, correct examples relative to hard
ones. It does not inspect class frequency at all, so a difficult example from
any class retains a large contribution --- though under skew the examples it
up-weights are disproportionately from the minority classes. The separate
$\alpha_t$ multiplier re-weights by \textbf{class frequency}. Zhao et al.
treat $\gamma$ as a constant and $\alpha$ as dataset-specific, to be set by
cross-validation.

\textbf{Source.} Zhao et al.~\cite{zhaoS2021}, who bring focal loss into
micro-expression recognition ``to alleviate the inefficient model training
caused by the class imbalance in the MEs Datasets''. They adopt the
formulation developed for one-stage object detection, where the imbalance
being corrected is between foreground and background regions rather than
between emotion classes. Per the bibliography's stated policy, the original
detection paper carries no entry here and the formulation is attributed to the
reviewed paper that reports it.

\textbf{As implemented.} Partly. The difficulty term is used as written, with
$\gamma = 2$ fixed throughout, matching Zhao et al.'s reported exponent --- an
inherited setting, not a measured optimum for this corpus. The implementation
computes $p_t$ from the hard target label using log-softmax, then multiplies a
\emph{label-smoothed} cross-entropy term (\autoref{eq:label-smoothing}) by
$(1-p_t)^{\gamma}$, with the class weight applied afterwards and optionally.
The frequency term $\alpha_t$ is where this thesis departs: frequency
correction is assigned to the sampler instead, and when balanced sampling is
active a runtime guard disables the loss-side class weights, leaving focal
difficulty weighting in place. Conflating the two factors is precisely the
error recorded at \autoref{sec:imbalance-double-correction}.

% -----------------------------------------------------------------------------
\subsection{Inline quantities}
\label{eq:ch3-inline}
% -----------------------------------------------------------------------------

Three quantities in Chapter~3 are not set as display mathematics but carry the
same evidential weight.

\textbf{Flow magnitude, $\sqrt{u^2+v^2}$} (\autoref{sec:brightness-constancy}).
The scalar magnitude of the flow field, described by Zhao et
al.~\cite{zhaoS2021} alongside the component fields $u(x,y)$ and $v(x,y)$.
This thesis retains the two components as separate input channels rather than
reducing them to this magnitude, so the direction of motion is preserved.

\textbf{Effective frame rate, $6600/N$ fps} (\autoref{sec:tim-implications}).
Sampling 33 frames from a clip of $N$ frames recorded at 200~fps yields this
effective rate: \textbf{100~fps} at the median clip length $N=66$, which is
SMIC's recording rate and half what CASME~II was built to provide, and roughly
\textbf{52~fps} for the longest clip at $N=126$, below the original CASME
database's 60~fps~\cite{yan2014}. Eighty-one of the 156 clips are downsampled
by a factor of two or more. No reviewed paper reports this cost.

\textbf{True inter-frame interval, $N/(200 \times 32)$ seconds}
(\autoref{sec:tim-implications}). Magnification runs on the 33 sampled frames
at the nominal 200~fps, but those frames span the clip's full duration, so the
real interval between them is this rather than $1/200$. A band specified as
5--25~Hz therefore corresponds to roughly 2.5--12.5~Hz physically at the
median, and the discrepancy grows with clip length. The realised pass-band is
consequently clip-dependent --- neither the 5--25~Hz derived from the duration
rule of Bai et al.~\cite{bai2021} nor constant across the corpus. This is a
declared divergence; magnifying before subsampling is declared as further
work.


% =============================================================================
\section{Chapter 4 --- Methodology}
\label{sec:eq-ch4}
% =============================================================================

% -----------------------------------------------------------------------------
\subsection{Sinusoidal positional encoding}
\label{eq:posenc}
% -----------------------------------------------------------------------------

\textbf{Appears at} \autoref{sec:transformer-config}.

\[
	\mathrm{PE}(p, 2i) = \sin\!\left(\frac{p}{10000^{2i/96}}\right), \qquad \mathrm{PE}(p, 2i+1) = \cos\!\left(\frac{p}{10000^{2i/96}}\right).
\]

\textbf{Symbols.} $p$ is the position in the sequence, one of the 32 time
steps; $i$ indexes feature pairs within the embedding, so $2i$ and $2i+1$ are
the even and odd feature positions; 96 is $d_{\text{model}}$, the embedding
width, appearing here as a literal rather than a symbol because the
configuration fixes it.

\textbf{What it does.} It supplies the order information that self-attention
cannot recover on its own. Without it, attention is permutation-equivariant:
permuting the input tokens permutes the outputs correspondingly, so a shuffled
clip is indistinguishable from an ordered one (\autoref{sec:self-attention}).
Each feature pair oscillates at a different frequency, geometrically spaced by
the $10000^{2i/96}$ denominator, so the vector at each position is distinct
and positions at a fixed offset stand in a consistent linear relation.

\textbf{Source.} No corpus source. The encoding is Vaswani et al.'s, from the
original transformer, which the bibliography deliberately omits under the
policy described at the head of this appendix; Chapter~4 states it as the
standard form without attribution. The thesis reaches the architecture through
Zhang et al.'s SLSTT~\cite{zhang2022} and its Vision Transformer
foundation~\cite{dosovitskiy2021}, as traced at
\autoref{sec:transformer-review}. Note that ViT itself uses \emph{learned}
positional embeddings, not these fixed sinusoids.

\textbf{As implemented.} Yes, as written. It is a fixed buffer, not a learned
parameter, and contributes nothing to any parameter count. Two details of the
surrounding configuration are easily missed: the positional-encoding stage
applies its own \texttt{Dropout(p=0.1)} to the sequence \emph{after} the
encoding is added, which is a separate application from the encoder layers'
dropout at the same rate; and although the sinusoids could be evaluated at
arbitrary positions, the implementation uses a finite precomputed buffer, so
generalisation to unseen sequence lengths is neither exercised nor guaranteed.

% -----------------------------------------------------------------------------
\subsection{The mean-of-folds macro F1 ceiling}
\label{eq:fold-ceiling}
% -----------------------------------------------------------------------------

\textbf{Appears at} \autoref{sec:mean-of-folds-rejected}.

\[
	\frac{10 \cdot \tfrac13 + 8 \cdot \tfrac23 + 7 \cdot 1}{25} = \frac{3.333\ldots + 5.333\ldots + 7}{25} \approx 0.627.
\]

\textbf{Symbols.} The three numerator terms are the 25 leave-one-subject-out
folds grouped by how many of the three classes appear in each fold's held-out
ground truth: \textbf{10} folds contain one class, \textbf{8} contain two,
\textbf{7} contain all three. The fractions $\tfrac13$, $\tfrac23$ and $1$ are
the maximum macro F1 each group can attain. The denominator is the fold count.

\textbf{What it does.} It computes an upper bound, fixed by protocol
structure before any model is trained. A macro F1 computed \emph{within} a
fold is the unweighted mean of that fold's three per-class F1 scores, and a
class absent from that fold's ground truth contributes zero no matter what the
model predicts. A single-class fold is therefore capped at $1/3$ and a
two-class fold at $2/3$. Averaging across folds averages seventeen sub-unity
ceilings into every score, bounding the mean-of-folds macro F1 at
approximately \textbf{0.627} --- a property of which subjects were held out,
identical for all twelve configurations.

Pooling avoids the ceiling entirely: accumulating $TP$, $FP$ and $FN$ across
all 25 folds and all 156 clips before computing any per-class F1 means a class
locally absent from one fold no longer forces a zero into its contribution.
That is the construction of \autoref{eq:uf1}.

\textbf{Source.} This thesis's own arithmetic over the fold composition of
\autoref{fig:fold-composition}. It is not a published result.

\textbf{As implemented.} Not applicable --- this is an argument for rejecting
a metric, not a component. Its consequence is directly implemented: pooled
macro F1 is the primary reported metric, and the mean-of-folds quantity stored
under \texttt{macro\_f1} in the run artefacts is not. A metric with a hard
ceiling below 0.7 cannot serve as a headline result. The same fold composition
biases mean-of-folds \emph{accuracy} in the same direction without imposing a
hard ceiling; for the full configuration it exceeds pooled accuracy by
\textbf{6.3 points}.

% -----------------------------------------------------------------------------
\subsection{Inline quantities}
\label{eq:ch4-inline}
% -----------------------------------------------------------------------------

\textbf{Sampler weight, $1/n_c$} (\autoref{sec:sampler-config}). Every
training example is assigned this weight by the
\texttt{WeightedRandomSampler}, where $n_c$ is the number of training examples
of that example's class \emph{within the current fold's training split}.
Sampling proceeds with \texttt{replacement=True}, drawing as many indices per
epoch as the split contains, which makes each represented class equally likely
in expectation while leaving individual batches uneven. This changes class
exposure during training, not the number of distinct recorded expressions.
It is also what makes the class-weighting rule well founded rather than merely
convenient: because the skew correction genuinely happens here, disabling its
loss-side equivalent avoids double-counting rather than removing the
correction (\autoref{eq:focal}).

\textbf{Parameter counts} (\autoref{sec:parameter-counts}). Four components
combine to give any configuration's total: the convolutional backbone
(\texttt{ThreeStreamCNNBackbone}, all three streams) \textbf{14,544}; the
transformer encoder (\texttt{SLSTTTransformer}) \textbf{348,736};
\texttt{RawPatchEmbedding}, the no-CNN fallback, \textbf{4,704}; and the
classifier head \textbf{483}. Their sums give the four distinct totals of
\autoref{tab:parameter-totals}. Neither SimAM nor EVM appears as a row: SimAM
contributes zero parameters by construction (\autoref{eq:simam}), and
magnification is a preprocessing step rather than a network component. Of the
transformer's total, 192 belong to the final \texttt{LayerNorm(96)} applied
after its four encoder layers.


% =============================================================================
\section{Chapter 2 --- Background}
\label{sec:eq-ch2}
% =============================================================================

% -----------------------------------------------------------------------------
\subsection{Eulerian magnification: the small-motion approximation}
\label{eq:evm-taylor}
% -----------------------------------------------------------------------------

\textbf{Appears at} \autoref{sec:eulerian-idea}.

\[
    I(x,t)=f(x+\delta(t))
       \approx f(x)+\delta(t)\frac{\partial f}{\partial x}.
\]

\textbf{Symbols.} $f$ is the image intensity profile; $x$ is spatial
position; $\delta(t)$ is a small time-varying displacement; $I(x,t)$ is the
observed intensity at a \emph{fixed} location $x$ at time $t$.

\textbf{What it does.} It is the first-order Taylor expansion that licenses
the entire Eulerian approach. Eulerian magnification examines intensity
changes at fixed image locations rather than explicitly tracking moving
points, and this approximation is what connects the two: for sufficiently
small displacement, the intensity change observed at a fixed pixel is
proportional to the displacement times the local spatial gradient. Motion can
therefore be manipulated by manipulating intensity.

\textbf{Source.} Bai et al.~\cite{bai2021}, describing the amplitude-based
mechanism. The technique originates with Wu et al., who carry no bibliography
entry under the policy described above and are attributed through the reviewed
paper that reports them.

\textbf{As implemented.} Yes --- but note that the approximation holds only
for small $\delta(t)$. It does not imply that arbitrary intensity changes
represent facial movement, nor that larger amplification always improves
recognition.

% -----------------------------------------------------------------------------
\subsection{Eulerian magnification: the amplified signal}
\label{eq:evm-amplified}
% -----------------------------------------------------------------------------

\textbf{Appears at} \autoref{sec:eulerian-idea}.

\[
    \widetilde I(x,t)\approx
       f(x)+(1+\alpha)\delta(t)\frac{\partial f}{\partial x}.
\]

\textbf{Symbols.} As above, with $\alpha$ the magnification factor and
$\widetilde I$ the reconstructed signal.

\textbf{What it does.} Adding $\alpha$ times the band-passed motion-related
variation back to the original produces a signal that, under the same
first-order approximation, resembles what would have been observed had the
displacement itself been $(1+\alpha)$ times larger. The mechanism is entirely
photometric; nothing is tracked or warped.

\textbf{Source.} Bai et al.~\cite{bai2021}, as above.

\textbf{As implemented.} \textbf{Not in the temporal filter.} The magnitude
relation above is implemented, but the band-pass stage that isolates
$\delta(t)$ differs: Bai et al. describe a Butterworth filter, whereas this
implementation masks Fourier coefficients outside the selected band and
applies the inverse transform (\autoref{sec:temporal-filtering}). The spatial
decomposition uses a Laplacian pyramid with the separable blur kernel
$[1,4,6,4,1]/16$ and factor-two downsampling. A second and larger divergence
concerns \emph{when} magnification is applied: it runs after subsampling, which
makes the realised pass-band clip-dependent (\autoref{eq:ch3-inline}).

% -----------------------------------------------------------------------------
\subsection{Brightness constancy}
\label{eq:brightness}
% -----------------------------------------------------------------------------

\textbf{Appears at} \autoref{sec:optical-flow}.

\[
    I(x,y,t)=I(x+u\Delta t,y+v\Delta t,t+\Delta t).
\]

\textbf{Symbols.} $I$ is image intensity; $(x,y)$ is spatial position; $t$ is
time; $u$ and $v$ are the horizontal and vertical components of local
velocity, for the continuous derivation.

\textbf{What it does.} It states the assumption on which optical flow rests:
a moving point retains its intensity as it moves. Everything estimated as
``motion'' downstream is motion only insofar as this holds.

\textbf{Source.} OFF-ApexNet, Liong et al.~\cite{liong2019a}.

\textbf{As implemented.} Yes, as the assumption underlying the Farneb\"ack
estimator used here --- with two caveats the thesis states rather than
assumes away. In the discrete pipeline $u$ and $v$ denote displacement in
pixels between selected frames rather than physical velocity per second. And
illumination changes or large inter-frame displacement can violate the
approximation outright, which the clip-dependent effective frame rate
(\autoref{eq:ch3-inline}) makes a live concern for the longer clips.

% -----------------------------------------------------------------------------
\subsection{The optical flow constraint}
\label{eq:flow-constraint}
% -----------------------------------------------------------------------------

\textbf{Appears at} \autoref{sec:optical-flow}.

\[
    I_xu+I_yv+I_t=0.
\]

\textbf{Symbols.} $I_x$, $I_y$ and $I_t$ are the partial derivatives of
intensity with respect to $x$, $y$ and $t$; $u$ and $v$ are as above.

\textbf{What it does.} It is the first-order expansion of brightness
constancy, and it is \textbf{one equation in two unknowns}. This
underdetermination is the \textbf{aperture problem}: a local edge constrains
motion normal to the edge, along the intensity gradient, but leaves motion
parallel to the edge entirely undetermined. Every practical estimator must
therefore add something --- a neighbourhood assumption, a local motion model,
or explicit regularisation. Flow fields are consequently never pure
measurements; they carry the assumptions of whichever estimator produced them.

\textbf{Source.} Liong et al.~\cite{liong2019a}.

\textbf{As implemented.} Yes. The constraint is resolved by Farneb\"ack's
polynomial expansion, following the estimator choice reported by Zhao et
al.~\cite{zhaoS2021}, with OpenCV settings: pyramid scale 0.5, three levels,
window size 15, three iterations, polynomial neighbourhood 5, polynomial
smoothing 1.2, flags 0. Flow is computed for the 32 adjacent pairs among the
33 sampled frames.

% -----------------------------------------------------------------------------
\subsection{The symmetric displacement-gradient tensor}
\label{eq:strain-tensor}
% -----------------------------------------------------------------------------

\textbf{Appears at} \autoref{sec:optical-strain}. Its normal and shear
components are also named at \autoref{sec:strain-tensor} in Chapter~3.

\[
    \varepsilon=\tfrac12\left(\nabla\mathbf{u}
                         +(\nabla\mathbf{u})^{\mathsf T}\right),
    \qquad
    \varepsilon_{xx}=\frac{\partial u}{\partial x},\quad
    \varepsilon_{yy}=\frac{\partial v}{\partial y},
\]
\[
    \varepsilon_{xy}=\varepsilon_{yx}
       =\tfrac12\left(\frac{\partial u}{\partial y}
                    +\frac{\partial v}{\partial x}\right).
\]

\textbf{Symbols.} $\mathbf{u}=[u,v]^{\mathsf T}$ is the displacement field;
$\nabla\mathbf{u}$ is its gradient; the superscript $\mathsf{T}$ denotes
transpose. Symmetrising the gradient makes $\varepsilon_{xy}=\varepsilon_{yx}$
by construction.

\textbf{What it does.} It converts displacement into \emph{deformation}.
Flow describes how much things moved; strain describes how neighbouring
displacements \emph{differ}. The normal components describe local stretching or
compression along each axis, and the off-diagonal terms describe shear. The
consequence that motivates using it at all: a spatially uniform translation
has zero displacement gradient and therefore zero strain, so rigid head
movement is suppressed while differential skin motion survives.

\textbf{Source.} The small-strain approximation used for optical strain,
Shreve et al.~\cite{shreve2011}.

\textbf{As implemented.} Yes, with the discretisation as a stated modelling
choice rather than a neutral detail, since it sets the scale at which local
variation is measured. Shreve et al. approximate the spatial derivatives by
central differences with offsets of approximately two to three pixels; Liong
et al.~\cite{liong2014a} use one-pixel offsets, and this implementation
follows the latter via NumPy's \texttt{gradient}. Two limits are declared
rather than argued away: zero strain under uniform translation does not imply
immunity to head \emph{rotation}, projection effects or flow-estimation error,
and differentiation may amplify noise. The reduction of this tensor to a
scalar is where the thesis diverges from the literature --- see
\autoref{eq:strain-implemented}.

% -----------------------------------------------------------------------------
\subsection{Convolution parameter count}
\label{eq:conv-params}
% -----------------------------------------------------------------------------

\textbf{Appears at} \autoref{sec:convolution}.

\[
    N_{\mathrm{params}}=C_{\mathrm{out}}
       (C_{\mathrm{in}}k_Hk_W+1).
\]

\textbf{Symbols.} $C_{\mathrm{in}}$ and $C_{\mathrm{out}}$ are the input and
output channel counts; $k_H$ and $k_W$ are the kernel height and width; the
$+1$ is one bias per output channel.

\textbf{What it does.} It makes weight sharing quantitative. Each output
channel's kernel spans all input channels over its local window, and the same
kernel is applied at every position --- so the parameter count is
\textbf{independent of image dimensions}. That independence is the point of
convolution: no separate detector is learned for each image location. What
does scale with image size is computation and activation memory, since both
grow with the number of positions processed.

\textbf{Source.} Standard; no corpus source and no citation.

\textbf{As implemented.} Yes. The implemented streams use $(1,3,3)$ kernels
with $(0,1,1)$ padding followed by $(1,2,2)$ max-pooling, so the convolutions
preserve temporal length and learn purely spatial filters while pooling halves
height and width. Because $k_T = 1$, \textbf{no temporal convolution is
performed} --- the use of \texttt{Conv3d} does not by itself establish
otherwise. The complete stream can still depend on multiple frames through
BatchNorm during training and through SimAM whenever it is enabled, both of
which use statistics spanning time.

% -----------------------------------------------------------------------------
\subsection{Batch normalisation}
\label{eq:batchnorm}
% -----------------------------------------------------------------------------

\textbf{Appears at} \autoref{sec:normalisation}.

\[
    \widehat x=\frac{x-\mu_c}{\sqrt{\sigma_c^2+\epsilon}},
    \qquad y=\gamma_c\widehat x+\beta_c.
\]

\textbf{Symbols.} $\mu_c$ and $\sigma_c^2$ are the mean and variance for
channel $c$; $\epsilon$ is a small constant preventing division by zero;
$\gamma_c$ and $\beta_c$ are the learned per-channel scale and shift;
$\widehat x$ is the standardised activation and $y$ the output.

\textbf{What it does.} It rescales each channel's activations to a
standardised distribution and then lets the network learn to undo or reshape
that standardisation through $\gamma_c$ and $\beta_c$ --- so normalisation is
imposed but not forced to be retained.

\textbf{Source.} \textbf{None.} Batch normalisation is Ioffe and Szegedy's,
but no paper in this thesis's corpus introduces or examines it, so it is
presented as standard practice and carries no citation. Adding one would
require a new bibliography entry, which the citation policy of this thesis
excludes.

\textbf{As implemented.} Yes. \texttt{BatchNorm3d} computes training
statistics over $(B,T,H,W)$ for each channel, and evaluation uses running
estimates accumulated during training. These activation statistics are a
distinct thing from the per-sample input standardisation performed by the
dataset loader, and the two should not be conflated. Related regularisers in
the same configuration: \texttt{Dropout3d} removes entire feature channels for
a sample rather than masking individual spatial values, and layer
normalisation --- which standardises a token's feature dimension independently
of other samples --- appears inside the transformer and in the classifier
head.

% -----------------------------------------------------------------------------
\subsection{Scaled dot-product attention}
\label{eq:attention}
% -----------------------------------------------------------------------------

\textbf{Appears at} \autoref{sec:self-attention}.

\[
    \operatorname{Attention}(Q,K,V)
      =\operatorname{softmax}\left(\frac{QK^{\mathsf T}}{\sqrt{d_k}}\right)V.
\]

\textbf{Symbols.} $Q$, $K$ and $V$ are the query, key and value matrices,
each a projection of the input tokens; $d_k$ is the key dimension; the
softmax is taken row-wise.

\textbf{What it does.} Queries and keys determine \emph{which} values
contribute to each output, and by how much: the dot product $QK^{\mathsf T}$
scores every token pair, the softmax turns those scores into weights summing
to one, and the result mixes the values accordingly. The $\sqrt{d_k}$ divisor
moderates dot-product magnitude as the key dimension grows, keeping the
softmax out of its saturated regime. Multiple heads learn projections into
different subspaces, whose outputs are concatenated and linearly mixed.

The property that matters most here is negative: without position
information, attention is \textbf{permutation-equivariant}. Permuting the
input tokens permutes the outputs correspondingly, so the operation cannot
identify temporal order on its own. That is what
\autoref{eq:posenc} exists to supply.

\textbf{Source.} Vaswani et al.'s transformer, attributed through Dosovitskiy
et al.~\cite{dosovitskiy2021}, whose Vision Transformer is the foundation this
thesis reaches the architecture through, together with Zhang et al.'s
SLSTT~\cite{zhang2022} (\autoref{sec:transformer-review}). The original
transformer paper carries no bibliography entry under the stated policy.

\textbf{As implemented.} Yes, within \texttt{nn.TransformerEncoder}, but over
a different axis than in the published SLSTT. Here attention operates over
\textbf{32 time steps} of 96-dimensional features, not over spatial patches:
ViT splits an image into patches and relates them spatially, whereas the
sequence here is produced by collapsing both spatial dimensions with
\texttt{AdaptiveAvgPool3d((32,1,1))}. The configuration is $d_{\text{model}} =
96$, $n_{\text{head}} = 8$, four encoder layers, feed-forward dimension 256,
dropout 0.1, pre-norm ordering, GELU activation.

% -----------------------------------------------------------------------------
\subsection{Label smoothing}
\label{eq:label-smoothing}
% -----------------------------------------------------------------------------

\textbf{Appears at} \autoref{sec:label-smoothing}.

\[
    L_{\mathrm{smooth}}
      =-(1-\epsilon)\log p_t
        -\frac{\epsilon}{C}\sum_{c=1}^{C}\log p_c.
\]

\textbf{Symbols.} $\epsilon$ is the smoothing strength; $C$ is the number of
classes; $p_t$ is the predicted probability of the true class and $p_c$ of
class $c$. The target it corresponds to is
$q_c=(1-\epsilon)\mathbf{1}[c=t]+\epsilon/C$, the one-hot vector blended with
uniform mass.

\textbf{What it does.} It removes the incentive to drive the true-class
probability to exactly one. Because a share $\epsilon$ of the target mass sits
on every class, the loss is minimised by a confident but not saturated
prediction. The effect is on \emph{calibration}, not on class balance: it
discourages extreme confidence uniformly and does not preferentially weight a
rare class.

\textbf{Source.} Named, but not examined, by Dosovitskiy et
al.~\cite{dosovitskiy2021}, who list it among the regularisation parameters
they tune. \textbf{No reviewed micro-expression paper uses it, discusses it,
or reports what it does under class skew.}

\textbf{As implemented.} Yes, with $\epsilon = 0.05$, applied \emph{before}
focal modulation --- the focal factor of \autoref{eq:focal} still uses the
unsmoothed true-class probability. Neither the smoothing strength nor the
focal exponent is varied across the architectural comparisons. Because the
corpus supports the choice neither way, it is declared as a choice rather than
defended as a finding, and it appears as such in the divergence ledger.

% -----------------------------------------------------------------------------
\subsection{Precision, recall and per-class F1}
\label{eq:prf}
% -----------------------------------------------------------------------------

\textbf{Appears at} \autoref{sec:confusion-matrix}. The $F1_c$ term also
appears within \autoref{eq:uf1}.

\[
    P_c=\frac{TP_c}{TP_c+FP_c},\qquad
    R_c=\frac{TP_c}{TP_c+FN_c},\qquad
    F1_c=\frac{2TP_c}{2TP_c+FP_c+FN_c}.
\]

\textbf{Symbols.} In a confusion matrix, entry $(i,j)$ counts examples whose
true class is $i$ and predicted class is $j$; the diagonal holds correct
predictions. For class $c$: $TP_c$ counts correct predictions of $c$, $FP_c$
counts incorrect predictions of $c$, and $FN_c$ counts missed examples of $c$.

\textbf{What it does.} Precision penalises false alarms and recall penalises
misses --- each is trivially maximised at the other's expense, which is why
neither is reported alone. F1 is their harmonic mean, so a high value requires
both to be reasonably high; the form above is the algebraic simplification
that avoids computing $P_c$ and $R_c$ separately.

\textbf{Source.} Standard; no corpus source and no citation.

\textbf{As implemented.} Yes, with one convention made explicit because it is
load-bearing under this fold structure: a degenerate denominator, which occurs
whenever a class is absent from a fold, is assigned \textbf{zero} rather than
the class being omitted from the score. That convention is exactly what
produces the ceiling of \autoref{eq:fold-ceiling}, and exactly what pooling
across folds avoids.

% -----------------------------------------------------------------------------
\subsection{Macro-averaged F1}
\label{eq:macro-f1}
% -----------------------------------------------------------------------------

\textbf{Appears at} \autoref{sec:macro-micro}.

\[
    F1_{\mathrm{macro}}=\frac1C\sum_{c=1}^{C}F1_c.
\]

\textbf{Symbols.} $C$ is the number of classes and $F1_c$ is the per-class F1
of \autoref{eq:prf}.

\textbf{What it does.} It gives every class equal weight regardless of how
many examples it has, which is the whole reason for preferring it here.
Contrast micro F1, which first sums $TP$, $FP$ and $FN$ \emph{across classes}:
in single-label multiclass classification each error supplies one false
positive and one false negative, so micro F1 equals accuracy exactly, and can
therefore be dominated by a frequent class. Macro F1 and accuracy describe
different aspects of performance, and reporting both does not remove the need
to inspect per-class errors directly.

\textbf{Source.} Standard. See et al.~\cite{see2019} define the \emph{pooled}
variant, MEGC's UF1, which is the construction this thesis uses
(\autoref{eq:uf1}).

\textbf{As implemented.} Yes, and it is the primary metric --- but only in
its pooled form. Pooling the confusion matrices across folds before computing
any per-class F1 is a genuinely different estimator from computing macro F1
within each fold and averaging the results (\autoref{sec:pooled-vs-fold}). The
distinction is the subject of \autoref{eq:fold-ceiling}, and no run artefact
stores the pooled quantity under any key: it is derived afterwards by the
reporting scripts under \texttt{tools/}.
